% 共通の組版設定。本文と証明は各論文の節ファイルに一度だけ置く。
\documentclass[11pt,a4paper]{ltjsarticle}
\usepackage[margin=24mm]{geometry}
\usepackage{amsmath,amssymb,amsthm,mathtools,booktabs,longtable}
\usepackage{mathrsfs}
\usepackage[haranoaji]{luatexja-preset}
\usepackage[unicode,hidelinks]{hyperref}
\setlength{\emergencystretch}{3em}
\hypersetup{bookmarksnumbered=true}
\pdfstringdefDisableCommands{\def\qgr{qgr}\def\HH{HH}\def\Perf{Perf}}
\providecommand{\PublicationVersion}{publication-v32}
\newtheorem{theorem}{定理}[section]
\newtheorem{proposition}[theorem]{命題}
\newtheorem{lemma}[theorem]{補題}
\newtheorem{corollary}[theorem]{系}
\theoremstyle{definition}\newtheorem{definition}[theorem]{定義}
\newtheorem{assumption}[theorem]{入力・仮定}
\theoremstyle{remark}\newtheorem{remark}[theorem]{注意}
\renewcommand{\proofname}{証明}
\newcommand{\HH}{\operatorname{HH}}
\newcommand{\qgr}{\operatorname{qgr}}
\newcommand{\coh}{\operatorname{coh}}
\newcommand{\Perf}{\operatorname{Perf}}
\newcommand{\Ext}{\operatorname{Ext}}
\newcommand{\Hom}{\operatorname{Hom}}
\newcommand{\Spec}{\operatorname{Spec}}
\newcommand{\Spf}{\operatorname{Spf}}
\newcommand{\GL}{\operatorname{GL}}
\newcommand{\gldim}{\operatorname{gldim}}
\newcommand{\OO}{\mathcal O}
\renewcommand{\AA}{\mathbb A}
\newcommand{\PP}{\mathbb P}

% 各原稿の引用だけを共通文献データから組版する。
\makeatletter
\newcommand{\UsePublicationReferences}[1]{%
  \@for\pub@reference:=#1\do{%
    \expandafter\def\csname pubref@\pub@reference\endcsname{}%
  }%
}
\makeatother

\title{三次曲面のPell力学と境界特異点\\\large 局所圏のmodulusとMilnor被覆のPrym偏極}
\author{}
\date{2026年11月18日\quad publication-v32}
\begin{document}
\maketitle
\begin{abstract}
三角形を無限遠境界に持つ滑らかなアフィン三次曲面上で、
四つのVieta involutionの合成によるconicの反復逆像を調べる。
反復曲線のアフィン部分は全て$\mathbb G_m$と同型であるが、
射影閉包は一つの境界点に二つのcuspを持つ。
そのPuiseux対はPell行列の係数$(b_n,c_n)$となり、
算術種数、局所Milnor数、対数標準閾値を明示的に与える。
アフィンperfect圏の通常のentropyが零となる一方、
射の極次数、正規化層、正規化欠損層は異なる増大率を持つ。
境界の先頭係数$K=\kappa^2$は解析的不変量であり、
重み付き主部の二周期matrix factorization圏をPell指数とともに分類する。
さらに主部の射影Milnor曲線の二次被覆から楕円曲線への有限写像を構成し、
そのJacobian引戻しの核を標識二捩れ、Prym内の楕円因子へのtheta制限次数を$b_nc_n$と決定する。
この構成を共通のCiani四次曲線と極端有理楕円曲面$X_{411}$に結び付ける。
実際の境界germには主部とTjurina数が異なる例があるため、
実germ、重み付き模型、そのMilnor被覆の各結論を区別する。
\end{abstract}
\tableofcontents

\section{問題と主結果}
本稿の対象は、複素数体上の滑らかなアフィン三次曲面
\begin{equation}\label{eq:boundary-paper-cubic}
 U=\{uvw-\tfrac12(u^2+v^2+w^2)-Au-Bv-Cw=h\}
\end{equation}
と、その通常の三次射影閉包$S$である。
無限遠の三角形を$D=D_u+D_v+D_w$とする。
曲面方程式を一つの変数について二次方程式と読むと、Vieta involution
\[
 \iota_u(u,v,w)=(2vw-u-2A,v,w)
\]
およびその巡回二式を得る。
合成は右から行い、以下では
\begin{equation}\label{eq:boundary-paper-word}
 g=\iota_v\iota_w\iota_u\iota_w,\qquad
 M=\begin{pmatrix}1&-2\\-2&5\end{pmatrix},\qquad
 M^n=\begin{pmatrix}a_n&-b_n\\-b_n&c_n\end{pmatrix},
 \quad \Lambda=3+2\sqrt2
\end{equation}
を固定する。この曲面族はFricke指標多様体の族でもある。
量子代数の中心として現れる背景と周期の比較は関連論文
\cite{CompanionFricke}で扱うが、以下の境界計算は式
\eqref{eq:boundary-paper-cubic}から出発し、非可換モジュライの存在定理を仮定しない。

正の次数成長を持つ自己同型についても、アフィン曲面の構造層は引戻しで固定される。
また一つの滑らかなconicを反復しても、アフィン曲線の同型類は変わらない。
そこで、成長が射の極次数と射影閉包のどの局所情報へ移るかを問う。
本稿では、同じPell列が境界付値、Puiseux対、局所圏の次元、
Prymに埋め込まれる楕円曲線の偏極次数を支配することを明示する。

\subsection{反復conicの境界と、局所化前後の成長}
$\kappa\in\mathbb C$について
\begin{equation}\label{eq:boundary-paper-admissible}
 \kappa^2\ne1,\qquad
 (\kappa^2+2B\kappa+2h)(\kappa^2-1)+A^2+C^2+2\kappa AC\ne0
\end{equation}
を仮定し、$C_n=g^{-n}\{v=\kappa\}\subset U$と置く。
この条件の下で全$C_n$は$\mathbb G_m$と同型である。
被約閉包を$\overline C_n$、その正規化を$\nu_n:\mathbb P^1\to\overline C_n$と書く。

\begin{theorem}[境界二枝と三つの測定値]\label{thm:boundary-paper-growth}
$n\ge1$とする。$\overline C_n$の唯一の特異点は
$p_w=[0:0:1:0]$であり、そこにPuiseux対$(b_n,c_n)$を持つ二枝が交わる。
二枝の交叉重複度は$b_nc_n$で、重み$(b_n,c_n)$の主部は
\[
 F_{n,\kappa}=y^{2b_n}-2\kappa x^{c_n}y^{b_n}+x^{2c_n}.
\]
従って
\[
 \mu_n=(2b_n-1)(2c_n-1),\qquad
 \operatorname{lct}_{p_w}(S;\overline C_n)=\frac{b_n+c_n}{2b_nc_n}.
\]
$g^n$の座標環における最小次数の最大値を$d_n$とすると
\[
 d_n=b_n+c_n,\qquad d_0=1,\quad d_1=7,\quad
 d_{n+2}=6d_{n+1}-d_n.
\]
一方、$\Perf(U)$上のDHKK entropyは全$t\in\mathbb R$で零であり、
射の極filtrationの歪み率は$\log\Lambda$である。
$\Perf(S)$の任意のsplit generatorを$E$とすると、対象列
\[
 Q_n=\OO_{\overline C_n},\qquad B_n=\nu_{n,*}\OO_{\mathbb P^1},
 \qquad Z_n=B_n/Q_n
\]
の$\delta_0(E,-)$の指数増大率は順に
$2\log\Lambda,\log\Lambda,2\log\Lambda$となる。
\end{theorem}
\begin{proof}
次数については定理\ref{thm:pell-exact-degree}、
境界二枝については定理\ref{thm:two-boundary-Puiseux}、
局所不変量については定理\ref{thm:boundary-torus-link}と
定理\ref{thm:boundary-lct}で証明する。
三つの成長の比較は命題\ref{prop:affine-pullback-zero-entropy}、
定理\ref{thm:filtered-pell-entropy}、
定理\ref{thm:proper-versus-affine-complexity}、
定理\ref{thm:normalization-defect-growth}による。
\end{proof}
対象列$Q_n,B_n,Z_n$を一つの自己関手の反復像と同定したわけではない。
その増大率と、アフィン自己関手のentropyを別の量として扱う。

\subsection{位相型の内部の解析modulus}
\begin{theorem}[重み付き模型の圏的分類と実germの区別]
\label{thm:boundary-paper-moduli}
重み付き模型の二周期dg圏
$\mathcal Q_{n,\kappa}=\operatorname{MF}(\mathbb C[[x,y]],F_{n,\kappa})$は
\[
 \mathcal Q_{n,\kappa}\simeq_{\rm Morita}\mathcal Q_{m,\kappa'}
 \quad\Longleftrightarrow\quad n=m\ \text{かつ}\ \kappa^2=(\kappa')^2
\]
を満たす。固定した$n$ではlink、Milnor数、Tjurina数、対数標準閾値は
$\kappa$によらないが、$K=\kappa^2$はこれらの圏を区別する。

実際の反復境界germについても$K$は解析的不変量であるが、完全不変量とは主張しない。
$A=B=C=0,h=1,\kappa=2,n=1$の実germは
$\mu=27,\tau=24$を持ち、$\tau=27$の重み付き主部と接触非同値である。
その通常の整数次数付きdg特異点圏も、主部のものとMorita非同値である。
\end{theorem}
\begin{proof}
命題\ref{prop:boundary-kappa-analytic-invariant}と
定理\ref{thm:weighted-boundary-Morita-classification}が分類を与える。
実germの計算は命題\ref{prop:actual-boundary-tjurina-defect}、
整数次数付きdg圏の区別は系\ref{cor:actual-germ-Zdg-distinction}にある。
\end{proof}
二周期MFでは相対Hochschild環がJacobian環を与え、
通常の整数次数付きdg特異点圏では絶対的$\mathrm{HH}^0$がTjurina環を与える。
後者による実germの区別を、そのまま前者の分類へ置き換えることはしない。

\subsection{同じmodulusを記録するPrymの楕円因子}
$b=b_n,c=c_n$とし、$C$を$F_{n,\kappa}=1$の滑らかな射影模型、
$\widetilde C$を$a^2=x$を添加した二次被覆の滑らかな射影模型とする。
次の定理では$\deg\Theta|_E$は楕円曲線上の\emph{線束の次数}を意味する。
対応する偏極isogenyの次数はその平方である。

\begin{theorem}[Prymの二捩れ核と整数偏極]\label{thm:boundary-paper-prym}
滑らかな楕円曲線
\[
 D_\kappa:T^2=r^4-2\kappa r^2+1,\qquad
 E_\kappa:\eta^2=z(z^2-2\kappa z+1)
\]
に対して巡回被覆の分解
\[
 \widetilde C\xrightarrow{q'}D_\kappa\xrightarrow{\psi}E_\kappa,
 \quad
 q'(a,y)=\left(\frac{y^{b/2}}{a^c},\frac1{a^{2c}}\right),
 \quad \psi(r,T)=(r^2,rT)
\]
があり、次数は$\deg q'=bc$, $\deg\psi=2$である。
合成$q=\psi q'$のJacobian引戻しの核は
\[
 \ker q^*=\{0,[(0,0)-O]\}\subset J(E_\kappa)
\]
である。$q'^*:J(D_\kappa)\hookrightarrow J(\widetilde C)$の像は
$\operatorname{Prym}(\widetilde C/C)$に含まれ、
\[
 \deg\bigl(\Theta_{J(\widetilde C)}|_{q'^*J(D_\kappa)}\bigr)=bc.
\]
Prym全体へのtheta制限の偏極型は
$(1^{b-1},2^{2bc-b-c})$である。
\end{theorem}
\begin{proof}
写像の構成と平方類の障害は定理\ref{thm:boundary-kummer-prym}、
巡回性と分岐は定理\ref{thm:prym-cyclic-inertia}、
核と偏極は定理\ref{thm:prym-elliptic-theta-degree}で証明する。
\end{proof}
$K$は標識二捩れ付き楕円曲線の粗modulusであり、
その標識は上の自然な引戻しの核としても回収される。
固定した$\kappa$に対して楕円因子の同型類は変わらないが、
theta制限次数は$b_nc_n$となり、指数率$2\log\Lambda$で増大する。
全$n$に共通するCiani四次曲線を経由すると三楕円分解が得られ、
$\kappa$を動かす楕円族は既知の極端有理楕円曲面$X_{411}$となる。

\subsection{先行研究との関係と本稿の範囲}
Fricke三次曲面のVieta作用、境界の双有理力学、entropyは
Cantat--Loray\cite{CantatLorayCayley}により研究されている。
Jang\cite{JangBoundary}はMarkov曲面のtropical作用を
$\mathbb R^2/\{\pm1\}$上の線型作用で記述する。
本稿はこれらの一般理論を新規結果とはせず、指定した語$g$とconicの列に対し、
最小次数、二枝の先頭係数、Puiseux対、正規化欠損を同時に計算する。

圏的複雑度にはDHKKの定義とsmooth proper圏での比較を用いる
\cite{BoundaryDHKK}。
閾値の計算はNewton非退化特異点に対するHowaldの方法
\cite{BoundaryHowald}、局所Weinstein模型はCasalsの既知構成
\cite{BoundaryCasals}の適用である。
二枝の重み付き解析分類はC\^amara\cite{BoundaryCamara}の一般分類と整合する。
MFの分類にはDyckerhoffのHochschild環の計算
\cite{BoundaryDyckerhoff}とMather--Yauの復元定理
\cite{BoundaryMatherYau}を使い、実germの整数次数付きdg圏には
Hua--Keller\cite{BoundaryHuaKeller}を適用する。
これらの定理では、Hochschild\emph{次元}の一致だけで圏の同値を結論しない。

Prymのnormと偏極はMumford\cite{BoundaryMumford}の一般理論に従う。
Ciani四次の三楕円分解、楕円曲線の二同種、
$X_{411}$の分類\cite{BoundaryMirandaPersson}も既知である。
本稿の具体的な比較は、反復境界の同じ$K$から有限写像を構成し、
平方類の障害、引戻しの核、誘導偏極の次数まで計算する点にある。
これらのPell反復への適用の個別の既出性については、なお文献比較を要する。

以上の結果は、明示した仮定の下で本文中に証明を置く。
ただし本稿のPrym構成は重み付き主部のMilnor曲線についての構成であり、
高重み項を持つ実germへの移行、偏極を保つPrym全体の直積分解、
元の非可換曲面の自己同値、filteredかつ同変なHMSを主張するものではない。
研究ノートにはこれらの未解決点と、再現計算の範囲を記す。

\section{周期の核を捉える境界付値とPell力学}
\label{sec:boundary-pell-dynamics}
本節では三次曲面のVieta作用を境界の単項式付値で観測し、
正確な次数列、二つの無理な固有方向、必要な境界blow-upの数を求める。
基礎体は$\mathbb C$、曲面は滑らかな
\[
 U=\{uvw-\tfrac12(u^2+v^2+w^2)-Au-Bv-Cw=h\}
\]
とし、$S$を通常の三次射影閉包、$D=D_u+D_v+D_w$を無限遠の三角形とする。
$\mathsf H=[D]=c_1(\mathcal O_S(1))$、$\mathsf H^2=3$である。
以下の$\deg_U$は座標環の最小全次数で、射影交点次数とは区別する。

\subsection{単項式付値の境界と、三つの区分線型反射}
$D$の各節点で二つの局所境界座標に正の実重みを与え、
対応する単項式付値を取る。閉錐の軸は隣接境界成分の
divisorial valuationの正倍として貼り合わせる。
この空間を$\mathcal V_D$と書く。定数体上では付値は自明とする。
これは指定した対$(S,D)$の境界付値の部分空間であり、
全Berkovich空間や任意の退化のessential skeletonを同定する主張ではない。

\begin{proposition}[境界付値の三角錐表示]
\label{prop:boundary-valuation-cone}
$\nu\in\mathcal V_D$の座標関数の極次数
\[
 (a,b,c)=(-\nu(u),-\nu(v),-\nu(w))
\]
は
\[
 \mathcal C=\{(a,b,c)\in\mathbb R_{\ge0}^3\setminus\{0\}:
                  2\max(a,b,c)=a+b+c\}
\]
を動き、この写像は同相である。
さらに
\begin{equation}\label{eq:boundary-theta}
 \Theta:(\mathbb R^2\setminus\{0\})/\{\pm1\}\xrightarrow{\sim}\mathcal C,
 \qquad (r,t)\longmapsto(|r|,|t|,|r+t|)
\end{equation}
は区分整数線型な同相である。
有理方向は、境界節点を有限回blow-upして得られる境界素因子の付値の方向に一致する。
\end{proposition}
\begin{proof}
頂点$[u:v:w:t]=[1:0:0:0]$の局所座標を$V,Z,T$とすると、
曲面方程式は
\[
 VZ-\tfrac12T(1+V^2+Z^2)-T^2(A+BV+CZ)-hT^3=0.
\]
陰関数定理から$T=2VZ\varepsilon$、$\varepsilon(0,0)=1$なる単元を得る。
$\nu(V)=\alpha,\nu(Z)=\beta$なら
$u=1/T,v=V/T,w=Z/T$の極次数は
$(\alpha+\beta,\beta,\alpha)$である。
巡回する三頂点の錐を境界成分の軸で貼ると$\mathcal C$となる。
最大の座標が残り二つの和であることから式\eqref{eq:boundary-theta}の逆も一意に定まる。
局所重みが互いに素な正整数ならEuclidの互除法に従うcorner blow-upでその素因子を得る。
逆にcorner blow-upでは隣接二つの原始重みの和を取るので、全て有理方向である。
\end{proof}

\begin{theorem}[Vieta作用の境界線型化]
\label{thm:boundary-vieta-linearization}
Vieta involution
\[
 \iota_u=(2vw-u-2A,v,w),\quad
 \iota_v=(u,2uw-v-2B,w),\quad
 \iota_w=(u,v,2uv-w-2C)
\]
は$\mathcal V_D$を保存する。$\iota_*\nu(f)=\nu(f\circ\iota)$とすると、
極次数では
\begin{equation}\label{eq:boundary-tropical-vieta}
 T_u(a,b,c)=(2\max(b,c)-a,b,c)
\end{equation}
と巡回二式で作用する。$\Theta$の下では行列
\[
 R_u=\begin{pmatrix}-1&-2\\0&1\end{pmatrix},\quad
 R_v=\begin{pmatrix}1&0\\-2&-1\end{pmatrix},\quad
 R_w=\begin{pmatrix}1&0\\0&-1\end{pmatrix}
\]
による作用となり、係数$A,B,C,h$によらない。
\end{theorem}
\begin{proof}
曲面の関係を使うと、積表示
\begin{equation}\label{eq:boundary-vieta-product}
 u\,\iota_u(u)=v^2+w^2+2Bv+2Cw+2h
\end{equation}
を得る。$\max(b,c)>0$なので右辺の極次数は$2\max(b,c)$。
$b=c$の場合も、単項式付値の先頭環では二つの二乗項は異なるLaurent単項式であり、
和は消えない。従って式\eqref{eq:boundary-tropical-vieta}を得る。
多項式表示$2vw-u-2A$に単純なmax演算を施すと、
$a=b+c$の錐で先頭項の消去を見落とすため、積表示を使うことが重要である。

付値そのものの不変性も局所座標で確かめられる。
上の頂点で$\alpha>\beta$なら像は$w$頂点に中心を持ち、
その座標$u'/w,v/w$の先頭項はそれぞれ$Z,V/Z$の非零定数倍となる。
$\alpha<\beta$なら$v$頂点で$V,Z/V$となる。
いずれもunimodularな単項式座標変換である。
他の二頂点の内点では$2vw$の先頭項が消えず、同様に単項式変換となる。
等重みの壁と軸は付値空間における連続性で含まれ、軸では素因子のgeneric pointを取る。
よって三錐からなる$\mathcal V_D$が保存される。

最後に恒等式$|x+y|+|x-y|=2\max(|x|,|y|)$を使えば
$T_i\Theta=\Theta R_i$が従う。
例えば$T_u$の新しい第一座標は$|r+2t|$であり、
$R_u(r,t)=(-r-2t,t)$の第二・第三絶対値は$|t|,|r+t|$である。
\end{proof}

この区分線型模型の一般理論は既知である。
Jang\footnote{S. u. Jang, \emph{Tropical Dynamics of Markov Surfaces},
\href{https://arxiv.org/abs/2306.11357v2}{arXiv:2306.11357v2},
Lemma 4.8 and Theorem 4.10（取得できたPDFは2025年12月版）。
別版の題名\emph{Vieta Involutions on Tropical Markov Cubics}や節番号と混同しない。}
は非Archimedean体上のMarkov曲面のtropical作用を同じ
$\mathbb R^2/\{\pm1\}$模型で記述する。
本節はその一般理論を新規とせず、定数体に自明な付値と境界の局所座標で
本稿の周期の核にある一つの写像を直接調べる。

\subsection{最低次数の完全なPell漸化式}
\begin{theorem}[係数に依存しない厳密な次数列]
\label{thm:pell-exact-degree}
$g=\iota_v\iota_w\iota_u\iota_w$と置く（右から合成）。
正則関数$f\in\mathcal O(U)$について、$\deg_U f$を
$f$を表す$u,v,w$の多項式の全次数の最小値と定める。
\[
 d_n=\max\{\deg_U(u\circ g^n),\deg_U(v\circ g^n),\deg_U(w\circ g^n)\}
\]
とすると、全ての$n\ge0$で
\begin{equation}\label{eq:pell-degree-sequence}
 d_0=1,\quad d_1=7,\quad d_{n+2}=6d_{n+1}-d_n,
 \qquad(d_n)=(1,7,41,239,1393,8119,\ldots).
\end{equation}
$\Lambda=3+2\sqrt2$と置けば
\[
 d_n=\frac{(1+\sqrt2)\Lambda^n+(1-\sqrt2)\Lambda^{-n}}2.
\]
さらに$M=\left(\begin{smallmatrix}1&-2\\-2&5\end{smallmatrix}\right)$、
$M^n=\left(\begin{smallmatrix}a_n&-b_n\\-b_n&c_n\end{smallmatrix}\right)$と書くと
\[
 d_n=b_n+c_n,\qquad d_n^2-2c_n^2=-1,
 \qquad d_n+c_n\sqrt2=(1+\sqrt2)^{2n+1}.
\]
射影交点次数は別の量であり、$n\ge1$について
\begin{equation}\label{eq:projective-pell-degree}
 (g^n)^*\mathsf H
 =c_nD_u+b_nD_v+(b_n+c_n)D_w,
 \qquad \mathsf H\cdot(g^n)^*\mathsf H=2d_n
\end{equation}
となる。従って$\lambda_1(g)=\Lambda$がこの計算からも直接従う。
\end{theorem}
\begin{proof}
一回の合成を展開した最高全次数項は、順に
\[
 g_u=4uv^2+\text{lower terms},\quad
 g_v=64u^2v^5+\text{lower terms},\quad
 g_w=8uv^3+\text{lower terms}.
\]
$N=\left(\begin{smallmatrix}1&2\\2&5\end{smallmatrix}\right)$、
$(p_n,q_n)^{\mathsf t}=N^n(1,1)^{\mathsf t}$と置く。
$n\ge1$で$g^n$の三座標の最高項は$u,v$だけの単項式で、
次数は$(p_n,q_n,q_n-p_n)$であることを帰納的に示す。
実際、$p_n,q_n>0$で、次の$u$座標の項$4u_nv_n^2$は
競合する$v_nw_n$より高次であり、次の$w$座標の$8u_nv_n^3$も同様である。
最後の$v$座標は$2u_{n+1}w_{n+1}-v_n-2B$だから、
次数は$(p_n+2q_n,2p_n+5q_n,p_n+3q_n)$になる。
最高係数はそれぞれ$2^{\text{次数}-1}$で消えない。
この帰納法は$A,B,C,h$の特殊値に依存しない。

座標環の全次数フィルトレーションの随伴次数環は
$\mathbb C[u,v,w]/(uvw)$である。最高項は$u,v$のみの非零単項式なので、
曲面関係を使っても次数を下げられない。
従ってこれらは単なる多項式表示の次数でなく、最小次数$\deg_U$である。
$N^2-6N+I=0$から式\eqref{eq:pell-degree-sequence}と閉形式を得る。
$N=\operatorname{diag}(1,-1)M\operatorname{diag}(1,-1)$で、
$b_n+c_n=q_n$。$(d_{n+1},c_{n+1})=(3d_n+4c_n,2d_n+3c_n)$と
初期値$(1,1)$は、単数$3+2\sqrt2$の乗算そのものである。
これがPellの式を与える。

境界付値では$R_vR_wR_uR_w=-M$なので、$g_*$は$M$で表される。
三つの境界素因子$D_u,D_v,D_w$は、それぞれ
$(r,t)=(0,1),(1,0),(1,-1)$に対応する。
$n\ge1$では$0<a_n<b_n<c_n$であり、
$g^n$の三座標の極次数の最大値は各素因子に沿って
$c_n,b_n,b_n+c_n$となる。
一般のaffine hyperplaneの引戻しではこれらの先頭係数は消えないため、
その零因子の閉包は表示した極因子と線型同値である。
これは双有理写像による$\mathsf H$の引戻しであり、
$\mathsf H D_i=1$から式\eqref{eq:projective-pell-degree}を得る。
$n=0$では$\mathsf H^2=3$なので、後式を$n=0$へ適用してはならない。
\end{proof}

\subsection{無理付値の二方向と、必要なblow-upの数}
\begin{theorem}[二次無理数の固有付値と非正則化]
\label{thm:irrational-boundary-nonregularization}
$g_*$の射影化した境界$\mathcal V_D/\mathbb R_{>0}$上の固定点は二つだけで、
極次数を規格化すると
\[
 v_+=(1,1+\sqrt2,\sqrt2),\qquad
 v_-=(1,\sqrt2-1,\sqrt2)
\]
となる。$g_*v_+=\Lambda v_+$、$g_*v_-=\Lambda^{-1}v_-$。
射影作用はnorth--south型で、$v_-$以外の全ての方向は正反復で$v_+$へ収束する。
有理方向には周期点がない。
従って、$S$から境界上の点を有限回blow-upして得られるどの滑らかな曲面にも、
$g$は$U$上の作用と両立する正則自己同型として延長しない。
\end{theorem}
\begin{proof}
$M$の二固有線は
$\mathbb R(1,-1-\sqrt2)$と$\mathbb R(1,-1+\sqrt2)$である。
$\Theta$を適用して二つの表示を得る。
異なる固有値$\Lambda>1>\Lambda^{-1}>0$による実射影直線上の作用はnorth--south型。
$M^m$の固有線も同じ二つなので、有理線の周期性は不可能である。

有限blow-up後にも元の三辺のstrict transformは境界素因子として残る。
もし$g$がその曲面の自己同型へ延長すれば、$U$の補集合である有限個の境界成分を
置換するため、各strict transformの付値は周期的となる。
しかしその付値は元の$D_i$と同じ有理方向であり、前半に矛盾する。
この議論は各blow-up中心が節点でなく境界の滑らかな点でも有効である。
\end{proof}

この結論は「有限blow-upによりalgebraic stabilityを得る」こととは異なる。
Algebraic stabilityは正則自己同型への延長を意味しない。
また、ここで扱ったのは指定した$S$を支配する有限境界blow-upであり、
無関係な全コンパクト化や非代数的コンパクト化の分類は主張しない。

\begin{proposition}[一つの反復付値を抽出するのに必要な正確な回数]
\label{prop:pell-corner-depth}
$\nu_n=g_*^n\operatorname{ord}_{D_w}$とする。
$n\ge1$では$\nu_n$の中心は$v$頂点で、二つの局所境界座標での原始重みは
$(p_n,q_n-p_n)$（順序を除く）である。
その比の有限連分数は
\[
 \frac{q_n-p_n}{p_n}
 =[1;\underbrace{2,\ldots,2}_{2n-2\text{個}},3].
\]
従って、この付値を境界素因子として抽出するためのcorner blow-upの最小回数は
正確に$4n$である。例えば極次数の最初の三つは
\[
 (3,7,4),\quad(17,41,24),\quad(99,239,140)
\]
であり、極次数は指数増大する一方、一つの指定付値の抽出深さは線型に増える。
\end{proposition}
\begin{proof}
$\Theta(M^n(1,-1))=(p_n,q_n,q_n-p_n)$で、最大は第二成分である。
頂点での局所座標の重みはその差から求まり、unimodularityにより互いに素である。
$R_n=(q_n-p_n)/p_n$は
$R_1=4/3$、$R_{n+1}=(4+3R_n)/(3+2R_n)$を満たす。
$R_n=1+1/T_n$と書けば$T_{n+1}=2+1/(2+1/T_n)$であるため、
各反復で連分数へ二つの$2$が追加される。

corner blow-upで重み$(a,b)$の付値を追うと、$a>b$なら次の重みは$(a-b,b)$、
$b>a$なら$(a,b-a)$となる。終端の$(1,1)$ではさらに一回のblow-upで素因子が現れる。
これは付値の中心を毎回blow-upする必要があるため最小の過程でもある。
従って回数は連分数の部分商の和で、$1+2(2n-2)+3=4n$となる。
\end{proof}

\subsection{何が周期から失われ、どこで回復されるか}
$g$の$H_2(U,\mathbb Z)$作用が恒等であることは、
Goldman--Neumann\cite{BoundaryGoldmanNeumann}のTheorem 2から従う。
実際、$(x,y,z)=(2u,2v,2w)$とすると同論文の係数は
$(P,Q,R)=(-4A,-4B,-4C)$、fiberの値は$-8h-2$となる。
各Vieta involutionは$H_2$上で$-I$として作用するため、四つの積$g$は恒等に作用する。
各involutionが留数形式の符号を反転することも局所座標で直接確認でき、
$g$は留数形式とその全periodを保つ。
本節では同じ$g$について、境界の原始付値の無限軌道、Pell次数、無理な安定方向を得た。
従って、Hodge拡大が分裂する場合でもしない場合でも、
これらの境界力学は同じである。係数に依存する周期と、
この写像の係数非依存なtropical境界を分ける必要がある。

Cantat--Loray\footnote{S. Cantat, F. Loray,
\emph{Dynamics on Character Varieties and Malgrange irreducibility of Painlev\'e VI equation},
Ann. Inst. Fourier 59 (2009), 2927--2978,
\href{https://doi.org/10.5802/aif.2512}{doi:10.5802/aif.2512}, \S3.1;
著者長稿\emph{Holomorphic dynamics, Painlev\'e VI equation and character varieties},
\href{https://perso.univ-rennes1.fr/serge.cantat/Articles/painleve_long.pdf}{\S4.2, Remark 4.5}では
境界節点の反復blow-upと円周作用を記述する。版による節番号の違いに注意する。}
の境界力学・entropy定理、Jangのtropical線型化自体は先行結果である。
ここでは本稿の周期の核にある明示的な語に適用し、最小次数、射影交点次数、
個別付値の抽出回数を区別した正確な計算としてまとめた。
Pell単数$1+\sqrt2$とその連分数は、同じ二次元整数行列を介して現れるのであり、
新たな算術同値や量子化を主張するものではない。

この境界付値の作用と局所化による複雑度の変化を次節で比較する。
wrapped Fukaya圏への作用、量子代数の自己同型の成長率、
非Archimedean good reductionについては別の構成が必要である。

\section{局所化が消す複雑度、極のfiltration、境界に蓄積する種数}
\label{sec:localization-entropy}
前節のPell成長を圏に移す際、まずどの量を測るかを区別しなければならない。
本節では同じ自己同型について、通常のDHKK entropyは零、
射の極次数の歪みは$\log(3+2\sqrt2)$、
曲線の射影閉包の対象複雑度は$2\log(3+2\sqrt2)$となることを示す。
最後の量は\emph{対象の列}の増大率であり、射影圏の自己関手のentropyを
構成したという主張ではない。

$U$は滑らかな複素LG曲面、$R=\mathcal O(U)$、$j:U\hookrightarrow S$は
通常の三次コンパクト化、$D=D_u+D_v+D_w$とする。
$g=\iota_v\iota_w\iota_u\iota_w$、$\sigma=g^*:R\to R$、
$\Lambda=3+2\sqrt2$を固定し、前節の記号
\[
 M^n=\begin{pmatrix}a_n&-b_n\\-b_n&c_n\end{pmatrix},
 \qquad d_n=b_n+c_n
\]
を使う。本節は\emph{可換な中心曲面}の$\Perf(U)$を扱う。
関連論文で扱う非可換曲面の圏への$g$の持上げや、
選択したLiouville構造のexact symplectic自己同型としての実現は別問題である。

\subsection{アフィン引戻しの通常の圏的entropyは零}
\begin{proposition}[固定生成元と射への非自明な作用]
\label{prop:affine-pullback-zero-entropy}
非空のアフィン$k$-scheme $X=\Spec R$の自己射$f$について、
$\mathbf Lf^*: \Perf(X)\to\Perf(X)$のDHKK entropyは
\[
 h_t(\mathbf Lf^*)=0\qquad(t\in\mathbb R)
\]
である。特に本稿の$g$では$h_t(\mathbf Lg^*)=0$である。
しかし生成元$\OO_U$の自己準同型環に対する作用は$\sigma=g^*$であり、
対象を固定することから関手が恒等だとは結論できない。
\end{proposition}
\begin{proof}
$G=R$は$\Perf(R)$のsplit generatorで、$(\mathbf Lf^*)^nG\simeq G$。
DHKKの三角塔による複雑度\footnote{
G. Dimitrov, F. Haiden, L. Katzarkov, M. Kontsevich,
\emph{Dynamical systems and categories},
\href{https://arxiv.org/abs/1307.8418}{arXiv:1307.8418},
Definitions 2.1, 2.4, Proposition 2.2, Lemma 2.3.}
には一段の塔から$\delta_t(G,G)\le1$がある。
任意の点$p\in X$を取り、exact関手
$-\otimes_R^{\mathbf L}k(p):\Perf(R)\to\Perf(k(p))$を適用すると
\[
 1=\delta_t(k(p),k(p))\le\delta_t(R,R)
\]
となる。従って$\delta_t(G,(\mathbf Lf^*)^nG)=1$で、定義から結論する。
ここで非properな$\Perf(U)$に、proper圏用の有限次元Ext公式を流用していない。
\end{proof}

この命題はDHKKの定義の初等的帰結である。
前節の正の次数成長を持つ写像にも適用できるため、
\emph{通常の$\Perf(U)$のentropyにPell成長を期待する方法は成立しない}。
成長を検出するには、少なくとも射のfiltration等の追加データが必要である。

\subsection{局所化の射に残る極次数}
境界を定義するsectionを$s_D$とし、$m\ge0$について
\begin{equation}\label{eq:pole-endomorphism-filtration}
 F_mR=\{f\in R:\operatorname{div}_S(f)+mD\ge0\}
 =\frac{H^0(S,\OO_S(mD))}{s_D^m}
 \subset\operatorname{End}_{\Perf(U)}(\OO_U)
\end{equation}
と置く。右辺の割算はsectionを有理関数として読む意味で、ベクトル空間の商ではない。

\begin{lemma}[次数filtrationと局所化表示]
\label{lem:pole-degree-localization}
$F_mR$は$u,v,w$の全次数$m$以下の多項式の像に一致し、
\[
 \dim F_mR=\binom{m+3}{3}-\binom m3
 =\frac{3m^2+3m+2}{2}.
\]
また$R=\varinjlim_m\Hom_S(\OO_S,\OO_S(mD))$であり、
そのcolimitのfiltrationが式\eqref{eq:pole-endomorphism-filtration}である。
\end{lemma}
\begin{proof}
$D$は超平面切断である。三次hypersurfaceの完全列と
$H^1(\PP^3,\OO(m-3))=0$により、全$m\ge0$で
斉次多項式から$H^0(S,\OO_S(mD))$への制限は全射となる。
$t=1$で非斉次化すれば次数filtrationを得る。
三次の一本の関係によるHilbert値は上式となる。
任意の$U$上の正則関数は、normalな$S$上で境界にのみ有限位数の極を持つので
colimitの主張も従う。
\end{proof}

これはC\^ot\'{e}--Kartalのcategorical compactificationによる
pole-order/localization filtrationの具体例である。\footnote{
L. C\^ot\'{e}, Y. B. Kartal,
\emph{Categorical action filtrations via localization and the growth as a symplectic invariant},
\href{https://doi.org/10.1515/crelle-2023-0037}{doi:10.1515/crelle-2023-0037};
\href{https://arxiv.org/abs/2108.05938v3}{arXiv:2108.05938v3},
\S5.2.2, (5.9)--(5.10), Corollary 5.26, Theorem 5.27,
\S7.1.1, Corollaries 7.1--7.2.}
本稿でこれから測るのは、静的な関数$m\mapsto\dim F_mR$ではなく、
自己同型が\emph{何倍の極次数を必要とするか}である。

\begin{definition}[射filtrationの歪み率]
\label{def:morphism-filtration-distortion}
有限次元生成部分空間$V\subset R$、$1\in V$に対し
$F_m^VR=V^m$とし、
\[
 \ell_V(n)=\min\{r\ge1:\sigma^n(V)\subset V^r\},\qquad
 h_{\rm pole}(\sigma)=\lim_{n\to\infty}\frac1n\log\ell_V(n)
\]
と定める。この率を通常のDHKK entropyとは区別する。
\end{definition}

\begin{theorem}[Pell成長のfilteredな検出]
\label{thm:filtered-pell-entropy}
上の極限は存在し、有限次元生成部分空間$V$の選択に依存しない。
本稿の$V=\langle1,u,v,w\rangle$では
\[
 \ell_V(n)=d_n,\qquad h_{\rm pole}(\sigma)=\log\Lambda>0.
\]
また、各$n$について$d_n$は全ての$m\ge1$に同時に
$\sigma^n(F_mR)\subset F_{d_nm}R$を満たす最小の整数である。
$\sigma^n$を共役する代数同型と、その逆がfiltrationを線型の尺度変更まで保つなら、
共役後もこの歪みの指数率は変わらない。
\end{theorem}
\begin{proof}
積を取ると$\ell_V(n+m)\le\ell_V(n)\ell_V(m)$なので、Feketeの補題から極限が存在する。
別の有限次元生成空間$V'$に対して$V\subset(V')^a$、$V'\subset V^b$となる$a,b$がある。
すると
$\sigma^n(V')\subset V^{b\ell_V(n)}\subset(V')^{ab\ell_V(n)}$であり、
逆向きも同様。従って両$\ell$は$n$によらない定数倍で比較され、指数率は同じである。
$V=F_1R$の場合、$\ell_V(n)$は三座標の最小次数の最大値そのもので、
前節の定理\ref{thm:pell-exact-degree}から$d_n$に等しい。
乗法性が全$m$の上界を与え、$m=1$が最小性を与える。
共役と逆共役に同じ線型比較を適用すれば最後の主張を得る。
\end{proof}

\begin{corollary}[一様な加法的shiftでは成長を吸収できない]
\label{cor:no-additive-filtered-lift}
極filtrationと線型同値なfiltration$G_mR$に対して、ある定数$C$で
$\sigma(G_mR)\subset G_{m+C}R$が全$m$で成り立つことはない。
従ってこのfiltrationを保つ自己同型、または次数を保つRees環自己同型による持上げも存在しない。
\end{corollary}
\begin{proof}
そのような包含を反復すれば$\sigma^n(G_mR)\subset G_{m+nC}R$となる。
有限生成空間$V$を一つの$G_m$に入れ、両filtrationの線型比較を用いると
$\ell_V(n)=O(n)$となり、正の指数率に矛盾する。
Reesの次数を保つ持上げは$C=0$の場合である。
\end{proof}

同じ射の像$\sigma^n(F_1R)$は全$n$で次元4のままである。
$\dim F_{d_n}R\sim\tfrac32d_n^2$は、\emph{その像を含む極次数予算}の大きさであり、
射の像そのものの次元が増えたことを意味しない。

\subsection{開曲面では二項、閉包では指数的な対象複雑度}
$v$座標の一般の値$\kappa\in\mathbb C$を一つ選び、
$C_0=\{v=\kappa\}\subset U$が滑らかな既約affine conicとなるようにする。
$C_n=g^{-n}(C_0)$、$\overline C_n\subset S$をその被約閉包とする。
全$C_n$は$C_0$と同型なので、滑らかで正規化の種数は零である。

\begin{theorem}[境界に集中する算術種数]
\label{thm:boundary-curve-genus}
全$n\ge0$で
\begin{align*}
 [\overline C_n]&=P_n=c_nD_u+b_nD_v+d_nD_w,\\
 D\cdot\overline C_n&=2d_n,\qquad P_n^2=4b_nc_n,\\
 p_a(\overline C_n)&=1+2b_nc_n-d_n,
 \qquad \chi(\OO_{\overline C_n})=d_n-2b_nc_n.
\end{align*}
特異点は境界にのみあり、その局所$\delta$不変量の総和は
$1+2b_nc_n-d_n$である。
最初の四つの$(\deg, p_a)$は
$(2,0),(14,14),(82,656),(478,23422)$となる。
\end{theorem}
\begin{proof}
$\sigma^n(v)-\kappa$の境界極次数は、前節の$\Theta$模型から
$(c_n,b_n,d_n)$である。
実際、三境界の有理線$(0,1),(1,0),(1,-1)$に$M^n$を作用させた後、
$\Theta$の第二成分を取ればそれぞれ$c_n,b_n,b_n+c_n$となる。
定数$\kappa$は負の付値を消さない。$n=0$での$D_v$上の位数は零である。
従って主因子の零部分の閉包と極部分は線型同値で、最初の式を得る。
$D_i^2=-1$、$D_iD_j=1$ ($i\ne j$)から二つの交叉式が従う。
$K_S=-D$とadjunctionにより$p_a=1+(P_n^2-DP_n)/2$となる。
閉包は既約有理曲線で、$U$上では滑らかなので、正規化の種数公式が最後の主張を与える。
この結果は各境界特異点のPuiseux型まで分類したものではない。
\end{proof}

\begin{theorem}[局所化前後で異なる対象の成長率]
\label{thm:proper-versus-affine-complexity}
$E$を$\Perf(S)$の任意のsplit generator、$Q_n=\OO_{\overline C_n}$とする。
このとき
\begin{equation}\label{eq:proper-object-growth}
 \lim_{n\to\infty}\frac1n\log\delta_0(E,Q_n)=2\log\Lambda.
\end{equation}
一方、$\Perf(U)$では
\[
 j^*Q_n\simeq\operatorname{Cone}
   (R\xrightarrow{\sigma^n(v)-\kappa}R),\qquad
 \delta_0(R,j^*Q_n)\le2.
\]
従って局所化は、この自然な閉包の列の指数的複雑度を一様有界な複雑度へ落とす。
\end{theorem}
\begin{proof}
$\eta(T)=\sum_i\dim\Hom_S(\OO_S,T[i])$は三角形に関して劣加法的である。
従って$\eta(Q_n)\le\eta(E)\delta_0(E,Q_n)$。
一方$\eta(Q_n)\ge|\chi(Q_n)|=2b_nc_n-d_n$ ($n\ge1$)であり、
$b_n,c_n,d_n$はいずれも正定数倍の$\Lambda^n$に漸近する。
これが指数率$2\log\Lambda$の下界を与える。

上界にはsmooth proper圏におけるDHKKの比較不等式\footnote{
Dimitrov--Haiden--Katzarkov--Kontsevich、前掲論文、Theorem 2.6の証明、式(9)。
任意の対象$T$について、固定生成元からの$\delta_0$と
$\sum_i\dim\Hom(E,T[i])$を$T$によらない定数で比較する。
この部分だけでproper性を使う。}
を用いる。$E^\vee$を有限個の固定vector bundleからなる有界複体で表す。
任意のそのようなbundle $B$について
$0\to B(-P_n)\to B\to B|_{\overline C_n}\to0$を使う。
$P_n$は合計$2d_n$個の境界成分を順に足した有効因子である。
途中の因子$T$の三係数は$d_n$以下で、$|T\cdot D_i|=O(d_n)$。
$B|_{D_i}$は固定された次数の線束の直和なので
\[
 \sum_j h^j(D_i,B(-T)|_{D_i})=O(d_n).
\]
$0\to B(-T-D_i)\to B(-T)\to B(-T)|_{D_i}\to0$を$2d_n$回使うと
$\sum_jh^j(S,B(-P_n))=O(d_n^2)$となる。
有限個のbundleとhypercohomologyについて同じ評価を足し、
$\sum_i\dim\Hom(E,Q_n[i])=O(d_n^2)$を得る。
DHKK比較によって$\delta_0(E,Q_n)=O(d_n^2)$となり、下界と合わせて式を得る。

アフィン側では$\sigma^n(v)-\kappa$は整域$R$の非零元なので、
表示した二項複体が$R/(\sigma^n(v)-\kappa)$を解消する。
二個の生成元のshiftからなる塔があるため、複雑度は2以下である。
\end{proof}

ここで$Q_n$が一つのexact自己関手の反復像であることは示していない。
閉包を取る操作はexact関手ではなく、式\eqref{eq:proper-object-growth}を
$\Perf(S)$の自己同値のentropyだと呼ぶことはできない。
ただし、もしあるexact自己関手$\overline F$と対象$Q_0$が
$\overline F^nQ_0\simeq Q_n$を満たすなら、三角塔の劣乗法性から
$h_0(\overline F)\ge2\log\Lambda$が必要となる。
そのようなliftの存在は本節の結論ではない。

\subsection{Hochschild、周期巡回、wrapped側で必要な情報}
滑らかなアフィン曲面のHKR同型は
\[
 \mathrm{HH}_0(R)=R,\qquad
 \mathrm{HH}_2(R)=\Omega_R^2=R\omega
\]
を与える。$g^*\omega=\omega$なので、$F_mR$と$F_mR\omega$により
両者に入るfiltrationの歪み率は同じ$\log\Lambda$である。
一方、周期巡回homologyとde Rham cohomologyの標準比較\footnote{
B. Antieau, \emph{Periodic cyclic homology and derived de Rham cohomology},
Ann. K-Theory 4 (2019),
\href{https://doi.org/10.2140/akt.2019.4.505}{doi:10.2140/akt.2019.4.505},
\S3およびIntroductionの標数零の古典比較への参照。
ここではsmoothな$\mathbb C$-代数に対する古典的mixed-complex HKRを用いる。}
から
\[
 \mathrm{HP}_0(R)\cong H^0_{\rm dR}(U)\oplus H^2_{\rm dR}(U),\qquad
 \mathrm{HP}_1(R)=0
\]
を得る。Goldman--Neumann\cite{BoundaryGoldmanNeumann}, Theorems 1--2により、
$U$の被約ホモロジーは次数2だけにあり、$H_2(U,\mathbb Z)\cong\mathbb Z^5$である。
前節で説明した$g$の$H_2(U)$上の恒等性から、
$g^*$は前者の次元$6$の空間にも恒等に作用する。
従って\emph{filtration付きの無限次元Hochschild群では見え、
有限次元のperiodic cyclic群と通常の対象entropyでは見えない}という区別がある。

\begin{proposition}[wrapped比較に対する十分条件と限界]
\label{prop:conditional-filtered-mirror-growth}
$\Psi:\Perf(U)\simeq\mathcal C$をexact同値とし、
$L=\Psi(\OO_U)$、$\Phi=\Psi\circ\mathbf Lg^*\circ\Psi^{-1}$とする。
すると$h_t(\Phi)=0$である。
さらに$\operatorname{End}_{\mathcal C}(L)$に乗法的filtration$\mathcal A_\ell$があり、
同型$\theta:R\xrightarrow{\sim}\operatorname{End}_{\mathcal C}(L)$が
極filtrationと$\mathcal A$を相互に線型の尺度変更まで保ち、
$\Phi(L)\simeq L$の選択の下で射への作用を共役するなら、
この射の反復歪み率は$\log\Lambda$である。
ここで$\mathcal A_\ell$の各部分は有限次元とし、歪みは
$\Phi^n(\mathcal A_m)\subset\mathcal A_{Lm}$を全$m\ge1$で満たす
最小整数$L$の指数率として測る。
\end{proposition}
\begin{proof}
通常のentropyは同値による共役で不変なので最初の命題による。
後半は定理\ref{thm:filtered-pell-entropy}の両方向のfiltration比較を用いる。
\end{proof}

C\^ot\'{e}--Kartalの結果は、対のHMSの下でpole-orderとReeb長のgrowthを
比較する具体的な枠組みを与える。ただし、\emph{growth関数の次元の比較だけから、
上命題に必要な特定自己同型との同変性や一様なfiltration比較は従わない}。
本稿ではこの$g$に対するfiltered・equivariantなHMSをまだ構成していない。
また$g$がholomorphic symplectic形式を保つことだけで、
選んだLiouville構造のexact symplectomorphismになるとはしない。
次に必要なのは、境界のlocalizationに適合するmirrorの対象と射を指定し、
その作用が上の極次数をどのaction filtrationへ移すかを確かめることである。

\begin{remark}[再現計算]
\texttt{verify\_localization\_entropy.py}は全係数の一反復、
Hilbert値、交叉恒等式、最初の25個の整数行列・種数を厳密に検算する。
圏的複雑度の全$n$の評価とHodge/巡回比較は本文の証明と引用した定理による。
\end{remark}

\section{境界の二つのcusp、トーラス絡み目、正規化欠損}
\label{sec:boundary-torus-links}
前節で得た算術種数の増大を、個々の枝とその交叉へ分解する。
同じ反復曲線は、開曲面では全て一つの滑らかなconicと同型であるが、
閉包では一つの境界点に二つのcuspを集める。
この局所模型からトーラス絡み目、有限位数のMilnor monodromy、
正規化の欠損層、Chebyshev型の実morsificationへ接続する。
以下の一般的な平面曲線・Milnor・Weinsteinの理論は既知として使い、
本稿のPell反復に対する具体的な同定と区別する。

\subsection{二つの無限遠枝の完全な先頭式}
$U,S,D,g$および$a_n,b_n,c_n,d_n$は前二節のものとする。
$C_0=\{v=\kappa\}\subset U$について
\begin{equation}\label{eq:boundary-conic-condition}
 \kappa^2\ne1,\qquad
 N_v(\kappa):=(\kappa^2+2B\kappa+2h)(\kappa^2-1)
             +A^2+C^2+2\kappa AC\ne0
\end{equation}
を仮定する。これは残余射影conicが滑らかで、二つの異なる無限遠点を持つ条件である。
$C_n=g^{-n}(C_0)$、その被約閉包を$\overline C_n$と書く。
以下は全て$n\ge1$についての主張である。

\begin{theorem}[Pell数をPuiseux対に持つ二枝]
\label{thm:two-boundary-Puiseux}
$\overline C_n$の特異点は境界頂点
$p_w=[0:0:1:0]$ただ一つで、そこに二つの解析的枝がある。
各枝は唯一のPuiseux対$(b_n,c_n)$を持ち、
枝同士の交叉重複度は$b_nc_n$である。

より正確に、$\rho_\pm=\kappa\pm\sqrt{\kappa^2-1}$とし、
出発conicの二つの無限遠点で$s=1/u$を局所座標とする。
$p_w$の滑らかな曲面局所座標$x=U/W=u/w$, $y=V/W=v/w$では、
$\rho=\rho_\pm$に対応する枝のパラメータは
\begin{equation}\label{eq:exact-boundary-branch-leading}
 x=2^{-b_n}\rho^{-a_n}s^{b_n}(1+O(s)),\qquad
 y=2^{-c_n}\rho^{-b_n}s^{c_n}(1+O(s)).
\end{equation}
従って重み$\operatorname{wt}(x)=b_n$, $\operatorname{wt}(y)=c_n$での
局所定義式の最低重み部分は、非零定数倍を除いて
\begin{equation}\label{eq:boundary-Newton-face}
 y^{2b_n}-2\kappa x^{c_n}y^{b_n}+x^{2c_n}
 =(y^{b_n}-\rho_+x^{c_n})(y^{b_n}-\rho_-x^{c_n}).
\end{equation}
一般には局所定義式全体がこの重み付き斉次式に解析的同値だとは主張しない。
\end{theorem}
\begin{proof}
出発conicは
\[
 u^2+w^2-2\kappa uw+2Au+2Cw+\kappa^2+2B\kappa+2h=0
\]
である。条件\eqref{eq:boundary-conic-condition}の下でその正規化は$\PP^1$、
無限遠の二点では$u=s^{-1}$、$w=\rho s^{-1}+O(1)$、$v=\kappa$となる。
$g^{-1}=\iota_w\iota_u\iota_w\iota_v$をこの順に右から作用させると
一反復の先頭項は
\[
 (u_1,v_1,w_1)=
 (16\rho^2s^{-5},\ 2\rho s^{-2},\ 64\rho^3s^{-7})(1+O(s))
\]
となる。積が元の座標と定数項より真に高い極を持つため、
ここでは先頭項の消去は起こらない。

一般に$n\ge1$で
\begin{equation}\label{eq:inverse-Pell-branch-leading}
 u_n=2^{c_n-1}\rho^{b_n}s^{-c_n}(1+O(s)),\quad
 v_n=2^{b_n-1}\rho^{a_n}s^{-b_n}(1+O(s)),\quad
 w_n=2^{d_n-1}\rho^{a_n+b_n}s^{-d_n}(1+O(s)).
\end{equation}
これを次の四Vietaに代入すると極次数は
$(c,b,c+b)\mapsto(5c+2b,2c+b,7c+3b)$、
$\rho$の指数は
$(b,a,a+b)\mapsto(2a+5b,a+2b,3a+7b)$となる。
$c_n=a_n+2b_n$とPell行列の漸化式から
式\eqref{eq:inverse-Pell-branch-leading}が帰納的に従う。
各先頭係数は$2^{\text{極次数}-1}$なので、全係数$A,B,C,h$で消えない。

$d_n>c_n>b_n>0$より二つの無限遠点はともに$p_w$に移る。
この点で射影三次式の$T/W$微分は$-1/2$なので$x,y$は正則局所座標であり、
式\eqref{eq:inverse-Pell-branch-leading}を割れば
式\eqref{eq:exact-boundary-branch-leading}を得る。
正規化$\PP^1$からproperな$S$への写像は二点でも延長する。
開曲面上では同型$g^{-n}$を合成しているから、そこで特異点はない。
境界の逆像はこの二点だけなので特異点は$p_w$に限られる。

$\gcd(b_n,c_n)=1$、$b_n\ge2$だから各枝の唯一のPuiseux対は$(b_n,c_n)$である。
また$ a_nc_n-b_n^2=1$から
\[
 \lim_{s\to0}\frac{y^{b_n}}{x^{c_n}}=\rho.
\]
二つの値$\rho_+,\rho_-$は非零で異なるため、二枝の最低重み式は
式\eqref{eq:boundary-Newton-face}となる。
一方の枝のパラメータを他方の定義式へ代入すると最低次数が$b_nc_n$で非零となるため、
交叉重複度も求まる。
\end{proof}

\begin{corollary}[toric解消後に残る元の係数]
\label{cor:boundary-trace-recovered}
重み$(b_n,c_n)$のprimitive rayを含むtoric変更を行うと、
二つの枝はその例外曲線の滑らかなtorus部分で、
$z=y^{b_n}/x^{c_n}=\rho_\pm$において横断的に交わる。
従って枝のPuiseux次数が増大しても、この標識付き例外曲線上では
$z+z^{-1}=2\kappa$により出発conicの係数を回復できる。
\end{corollary}
\begin{proof}
primitive性により例外曲線のtorus部分は滑らかである。
最低重み式を例外因子で割った制限は$z^2-2\kappa z+1$となり、二根は単純である。
高重み項はそこで消えるので、横断性と表示を得る。
\end{proof}

\subsection{絡み目、conductor、有限位数monodromy}
以下では$b=b_n,c=c_n$、
$\delta_n=1+2bc-b-c$と略記する。

\begin{theorem}[二成分トーラス絡み目と局所不変量]
\label{thm:boundary-torus-link}
$p_w$における正の代数的絡み目は$T(2b,2c)$である。
各成分は$T(b,c)$、成分間のlinking numberは$bc$であり、
\[
 \delta(\text{各枝})=\frac{(b-1)(c-1)}2,\quad
 \delta(\overline C_n,p_w)=\delta_n,\quad
 \mu_n=(2b-1)(2c-1)=2\delta_n-1.
\]
正規化の完備局所環を
$\widetilde{\mathcal O}=\mathbb C[[s_+]]\oplus\mathbb C[[s_-]]$とすると、
conductor idealは
\begin{equation}\label{eq:boundary-conductor}
 \mathfrak c_n=s_+^{\delta_n}\mathbb C[[s_+]]
               \oplus s_-^{\delta_n}\mathbb C[[s_-]].
\end{equation}
Milnor fiberは二境界成分と種数$\delta_n-1$を持つ。
その$H_1$上のmonodromyの特性多項式は
\begin{equation}\label{eq:Pell-link-monodromy}
 \det(TI-h_{n,*})
 =\frac{(T^{2bc}-1)^2(T-1)}{(T^{2b}-1)(T^{2c}-1)}.
\end{equation}
このhomology上のmonodromyの位数は正確に$2bc$で、スペクトル半径は1である。
\end{theorem}
\begin{proof}
式\eqref{eq:boundary-Newton-face}は二つの単純な非零rootを持つNewton非退化面である。
各枝の一つのPuiseux対と相互交叉数が局所embedded topologyを決定する。
同じ面の係数を二根が衝突しないように動かして
$x^{2c}+y^{2b}$へ移せるので、絡み目は$T(2b,2c)$となる。
これは平面曲線のPuiseux--link対応の標準適用であり、
解析的な右同値の主張ではない。\footnote{
R. Casals, \emph{Lagrangian skeleta and plane curve singularities},
J. Fixed Point Theory Appl. 24 (2022), 34,
\href{https://doi.org/10.1007/s11784-022-00939-8}{doi:10.1007/s11784-022-00939-8},
\S2.1, Remark 2.1のPuiseuxと代数的絡み目の対応も参照。}
一枝の半群$\langle b,c\rangle$のgap数と、二枝の交叉重複度を足して$\delta_n$を得る。
平面曲線の$\mu=2\delta-r+1$で$\mu_n$が従う。

各枝のconductor指数は、その枝のMilnor数と他枝との交叉数の和であるから
$(b-1)(c-1)+bc=\delta_n$となる。\footnote{
D. Kerner, V. Vinnikov,
\emph{On the determinantal representations of singular plane curves},
\href{https://arxiv.org/abs/0906.3012v1}{arXiv:0906.3012v1},
\S2.2.3, Definition--Proposition 2.2, (7)を参照する。
取得した2009年6月版の番号を用いる。}
同じ公式は$dx/f_y$の各枝上の位数からも得られる：
$\operatorname{ord}f_y=2bc-c$、$\operatorname{ord}dx=b-1$だから差は$-\delta_n$である。
Milnor fiberの$H_1$のrankは$\mu_n$、境界は二成分なので種数は$\delta_n-1$。

絡み目のMilnor fibrationは重み付き斉次模型と同じである。
模型$x^{2c}+y^{2b}$では固有値は
$\exp(2\pi i(i/(2c)+j/(2b)))$
($1\le i<2c,1\le j<2b$)。これはMilnor--Orlikの重み公式の二変数の場合で、
そのdivisor積を展開すると式\eqref{eq:Pell-link-monodromy}を得る。\footnote{
Milnor--Orlikの公式の確認箇所として
C. Hertling, M. Mase,
\emph{The combinatorics of weight systems and characteristic polynomials of isolated quasihomogeneous singularities},
J. Algebr. Comb. 56 (2022), 929--954,
\href{https://doi.org/10.1007/s10801-022-01138-x}{doi:10.1007/s10801-022-01138-x},
Theorem 6.4(a)を用いる。原公式はJ. Milnor, P. Orlik, Topology 9 (1970), 385--393。}
重み作用からmonodromyは$2bc$乗して恒等となる。
一方$b$は偶数、$c$は奇数で互いに素だから$\gcd(b+c,2bc)=1$であり、
$i=j=1$の固有値はprimitiveな$2bc$乗根である。従って正確な位数も$2bc$となる。
\end{proof}

\begin{center}
\begin{tabular}{c|c|r|r|r}
$n$&絡み目&$\delta_n$&$\mu_n$&$h_{n,*}$の位数\\\hline
1&$T(4,10)$&14&27&20\\
2&$T(24,58)$&656&1311&696\\
3&$T(140,338)$&23422&46843&23660
\end{tabular}
\end{center}
各$n$の局所monodromyは有限位数でも、その位数の指数増大率は
$2\log\Lambda$である。これは曲面自己同型$g$の反復とは異なる作用であり、
局所Milnor monodromyに正のスペクトル半径を見いだしたという主張ではない。

\subsection{正規化が分離する二つの対象成長率}
$\nu_n:\PP^1\to\overline C_n\subset S$を正規化とし、$S$上のsheaf
\[
 Q_n=\mathcal O_{\overline C_n},\qquad
 B_n=\nu_{n,*}\mathcal O_{\PP^1},\qquad
 Z_n=B_n/Q_n
\]
を考える。$S$はsmoothなので、これらは全て$\Perf(S)$の対象である。

\begin{theorem}[一つの境界点が持つ正規化欠損の成長]
\label{thm:normalization-defect-growth}
$Z_n$は$p_w$だけに支持された長さ$\delta_n$のsheafであり、
$H^1(\overline C_n,\mathcal O_{\overline C_n})\cong H^0(S,Z_n)$である。
$E$を$\Perf(S)$の任意のsplit generatorとすると
\begin{equation}\label{eq:normalization-complexity-rates}
 \lim_{n\to\infty}\frac1n\log\delta_0(E,B_n)=\log\Lambda,
 \qquad
 \lim_{n\to\infty}\frac1n\log\delta_0(E,Z_n)=2\log\Lambda.
\end{equation}
前節の$Q_n$の率は$2\log\Lambda$である。
$j:U\hookrightarrow S$で制限すると
$j^*B_n=j^*Q_n=\mathcal O_{C_n}$は二項の自由解消を持ち、$j^*Z_n=0$となる。
これらは対象列の増大率であって、新しい自己関手のentropyではない。
\end{theorem}
\begin{proof}
正規化完全列$0\to Q_n\to B_n\to Z_n\to0$と前定理から
$\operatorname{length}Z_n=\delta_n$、支持は$p_w$と分かる。
$\overline C_n$は既約proper、正規化は$\PP^1$なので、
両者の$H^0(\mathcal O)$は$\mathbb C$で写像は同型、$H^1(\PP^1,\mathcal O)=0$。
cohomology完全列から$H^1(Q_n)=H^0(Z_n)$を得る。

$Z_n$のcomposition seriesは$\delta_n$個の$k(p_w)$を因子に持つ。
三角塔の劣加法性から$\delta_0(E,Z_n)=O(\delta_n)$。
逆向きは$\sum_i\dim\Hom(\mathcal O_S,Z_n[i])=\delta_n$と
同じ量の三角塔に関する劣加法性から従い、
$\delta_0(E,Z_n)$は$\delta_n$の正定数倍で上下から挟まれる。

$B_n$については$\nu_n^*\mathcal O_S(1)=\mathcal O_{\PP^1}(2d_n)$だから、
\[
 \sum_i\dim\Hom_S(\mathcal O_S(1),B_n[i])=2d_n-1.
\]
これが$\delta_0(E,B_n)$の定数倍の$d_n$による下界を与える。
上界には前節で用いたsmooth proper圏のDHKK比較を使う。
$E^\vee$を固定vector bundleの有限有界複体で表す。
その一つ$V$について、$V(mH)$と$V^\vee(mH)$がともに大域生成となる$m$を選ぶ。
$\nu_n^*V$の$\PP^1$上の全分裂次数は$[-2md_n,2md_n]$に入るため、
全cohomology次元は$O(d_n)$である。
有限個のbundleに同じ評価を適用して
$\sum_i\dim\Hom(E,B_n[i])=O(d_n)$、従って$\delta_0(E,B_n)=O(d_n)$となる。
$d_n\asymp\Lambda^n$、$\delta_n\asymp\Lambda^{2n}$から表示した極限を得る。
開曲面では正規化が同型なので、最後の主張は前節の二項解消による。
\end{proof}

\subsection{Chebyshev morsificationと局所Weinstein模型}
\begin{proposition}[実divideを伴う局所symplectic模型（既知構成の適用）]
\label{prop:Pell-Chebyshev-Weinstein}
各$n$の絡み目の重み付き斉次代表を$x^{2c}-y^{2b}$とする。
非零の係数の尺度変更を許せば、その実morsificationは
\begin{equation}\label{eq:Pell-Chebyshev-morsification}
 F_\epsilon(x,y)=\epsilon^{2bc}
 \left[T_{2c}(x/\epsilon^b)-T_{2b}(y/\epsilon^c)\right],\quad \epsilon>0
\end{equation}
で与えられる。臨界点は全て実Morseで、総数は$\mu_n$。
零level上の実nodeが$\delta_n$個、残りの極値点が$\delta_n-1$個である。

この絡み目のmaximal-tb Legendrian近似$\mathcal L_n$に対する
Casalsの構成により、Weinstein pair$(\mathbb C^2,\mathcal L_n)$は
種数$\delta_n-1$、境界二成分のexact Lagrangian曲面と
$\mu_n$個のLagrangian thimbleからなるclosed arboreal skeletonを持つ。
二つの境界をhandleで閉じるWeinstein四多様体はEuler数$3$を持つ。
\end{proposition}
\begin{proof}
$T_m(\cos\theta)=\cos(m\theta)$より、臨界点は
\[
 x=\epsilon^b\cos\frac{i\pi}{2c},\quad
 y=\epsilon^c\cos\frac{j\pi}{2b},\qquad
 1\le i<2c,\quad1\le j<2b
\]
で、いずれも非退化である。臨界値は
$\epsilon^{2bc}((-1)^i-(-1)^j)$。
$i,j$が同じparityなら値は零で、Hessianの符号は異なるため実node。
その数は$bc+(b-1)(c-1)=\delta_n$である。
異なるparityの個数は$c(b-1)+b(c-1)=\delta_n-1$。
$\epsilon\to0$で式\eqref{eq:Pell-Chebyshev-morsification}は
$2^{2c-1}x^{2c}-2^{2b-1}y^{2b}$に収束する。
全臨界点がMilnor球内でこのように分裂するので、実morsificationを得る。
Chebyshevによるこの構成自体は既知である。\footnote{
Casals、前掲論文、Example 1.3後のChebyshevの説明、
Theorem 1.1およびCorollary 1.2。
ここでは計算した局所絡み目にその定理を適用する。}
同定理のLagrangian曲面はMilnor fiberにsmoothly isotopicなので、
種数と境界数は前定理で決まる。Morse点一つに一つのthimbleを対応させる。
最後のWeinstein四多様体は四球へ二つの2-handleを付けるのでEuler数は$1+2=3$となる。
\end{proof}

これは境界特異点に付随する\emph{局所模型}の構成である。
元のLG曲面のEuler数は$6$なので、最後の四多様体をその全体と同一視できない。
また、異なる$n$の模型間のexact functor、$g$のwrapped自己同値、
極filtrationとの一様な同変比較を構成したわけではない。
しかし、前節まで抽象的だった「境界局所化のmirror対象」に対し、
絡み目・実divide・thimbleの個数を指定した具体的な比較先を与える。

\begin{remark}[何を証明し、何を検算したか]
二枝の先頭項と全$n$の帰納法から、Puiseux型と相互交叉を決めた。
局所link、conductor、Milnor公式、DHKK比較、Casalsの構成は既知理論の適用である。
\texttt{verify\_boundary\_torus\_links.py}は一反復の全係数先頭式、
帰納漸化式、24反復の整数式、$n=1$模型のJacobi次元27とconductor14、
monodromy多項式、Chebyshevの臨界値resultantを厳密に再現する。
有限回の計算を全$n$の証明とせず、解析的右同値やglobal mirror同値も主張しない。
個別のPell反復への適用の文献上の新規性は網羅確認未了である。
\end{remark}

\subsection{Newton閾値と境界matrix factorizationの非同値性}
\label{subsec:boundary-thresholds-mf}
前の定理で得た二枝のNewton面を、乗数イデアルと局所特異点圏に適用する。
$n\ge1$、$b=b_n,c=c_n$、$\Lambda=3+2\sqrt2$と略記し、
式\eqref{eq:boundary-conic-condition}を仮定する。
定理\ref{thm:two-boundary-Puiseux}より、唯一の境界特異点は$p_w$であり、
局所被約方程式$f_n$は重み$(b,c)$の最低部分
\[
 (y^b-\rho_+x^c)(y^b-\rho_-x^c)
 =y^{2b}-2\kappa x^cy^b+x^{2c}
\]
を持つ。$\rho_+\ne\rho_-$、$\rho_+\rho_-=1$なのでNewton非退化である。
ここでの不変量は\emph{閉包曲線の境界特異点}のものであり、
曲面$U$そのものや元の非可換H3圏の不変量ではない。

\begin{theorem}[局所対数標準閾値と乗数イデアル]
\label{thm:boundary-lct}
$0<\tau<1$について、$p_w$での乗数イデアルは
\[
 \mathcal J(S,\tau\overline C_n)_{p_w}
 =\left\langle x^iy^j:
 b(i+1)+c(j+1)>2bc\tau\right\rangle
 \subset\mathcal O_{S,p_w}^{\rm an}.
\]
特に
\begin{equation}\label{eq:pell-lct}
 \boxed{\operatorname{lct}_{p_w}(S;\overline C_n)
 =\frac{b+c}{2bc}},\qquad
 \lim_{n\to\infty}-\frac1n\log\operatorname{lct}_{p_w}(S;\overline C_n)
 =\log(3+2\sqrt2).
\end{equation}
ここでは境界$D$に係数1を与えておらず、対$(S,D)$に対する相対閾値ではない。
\end{theorem}
\begin{proof}
Newton非退化性により、前の二枝のNewton面に適合するトーリック解消は$f_n$のlog resolutionでもある。
正の原始ray$(r,s)$に対応する例外曲線では、log discrepancyは$r+s$、
$f_n$の位数は$\min(2cr,2bs)$である。
$\tau<1$ではstrict transformは乗数イデアルへ条件を加えず、
各単項式の条件は全rayについて
$(i+1)r+(j+1)s>\tau\min(2cr,2bs)$となる。
Newton多面体の唯一のコンパクト面に対応する$(r,s)=(b,c)$の条件が
他の条件を含意し、表示式を得る。
これは非退化主部に対するHowaldのトーリック計算法の局所版を直接適用したものでもある。\footnote{
J. Howald, \emph{Multiplier Ideals of Sufficiently General Polynomials},
\href{https://arxiv.org/abs/math/0303203}{arXiv:math/0303203}, Theorems 11--12 and Corollary 13.
本稿の局所冪級数では、上記のNewton面と明示的トーリック解消を使って計算している。}
定数$1$が初めて乗数イデアルから外れる値が$(b+c)/(2bc)$である。
$b_n,c_n$はともに正定数倍の$(3+2\sqrt2)^n$に漸近するので極限式も従う。
\end{proof}

\begin{corollary}[境界matrix factorizationのHochschild次元]
\label{cor:boundary-mf-growth}
$\mathcal M_n=\operatorname{MF}(\mathbb C[[x,y]],f_n)$を
上で選んだ局所被約方程式の$\mathbb Z/2$-graded dg圏とする。
このとき
\[
 \mathrm{HH}^{\bar0}(\mathcal M_n)
 \cong\frac{\mathbb C[[x,y]]}{(\partial_x f_n,\partial_y f_n)},
 \qquad\mathrm{HH}^{\bar1}(\mathcal M_n)=0,
\]
\[
 \dim\mathrm{HH}^{\bar0}(\mathcal M_n)
 =(2b_n-1)(2c_n-1),\qquad
 \lim_{n\to\infty}\frac1n\log\dim\mathrm{HH}^{\bar0}(\mathcal M_n)
 =2\log(3+2\sqrt2).
\]
特に異なる$n\ge1$の$\mathcal M_n$は$\mathbb Z/2$-graded dg Morita同値ではない。
\end{corollary}
\begin{proof}
$f_n$は孤立した臨界点を持ち、形式二変数環はregularである。
従ってDyckerhoffの既知の計算\footnote{
T. Dyckerhoff, \emph{Compact generators in categories of matrix factorizations},
\href{https://arxiv.org/abs/0904.4713v5}{arXiv:0904.4713v5},
Corollary 6.5（有限rankの二周期dg圏）。
Theorem 6.6はhomology側の対応であり、ここではcohomology側を使用する。}
を適用できる。そのJacobian環の次元は局所Milnor数なので、定理\ref{thm:boundary-torus-link}から上の公式を得る。
$b_n,c_n$はともに真に増大し、Hochschild cohomologyはdg Morita不変だから
異なる$n$の同値は存在しない。
\end{proof}

ここで$\mathcal M_n$は\emph{閉包曲線の境界特異点}に付随する圏である。
これを$\operatorname{Perf}(U)$、元の非可換H3圏、あるいは
$\operatorname{Perf}(S)\to\operatorname{Perf}(U)$の核と同一視したわけではない。
同じ曲線のアフィン部分は全て$\mathbb G_m$である一方、
境界の特異点圏はHochschild次元だけでも相異なる。
この圏の\emph{列}の次元成長も、単一の自己関手のentropyではない。

\begin{remark}[同じアフィン曲線と異なる境界の圏]
最初の三例では
\[
\begin{array}{c|c|c|c}
 n&(b_n,c_n)&\dim\mathrm{HH}^{\bar0}(\mathcal M_n)&\operatorname{lct}\\ \hline
 1&(2,5)&27&7/20\\
 2&(12,29)&1311&41/696\\
 3&(70,169)&46843&239/23660
\end{array}
\]
となる。アフィン部分$C_n$は全て$\mathbb G_m$で、
通常のアフィン対象複雑度は有界であるにもかかわらず、
境界のMF圏はHochschild次元だけで互いに区別できる。
一方、局所Milnor monodromyは各$n$で有限位数である。
これらの異なる意味の成長を単一の自己関手のentropyにまとめない。

HowaldとDyckerhoffの一般理論は既知である。
本稿では指定したPell反復族に適用し、閾値の減衰率$\log\Lambda$と
局所圏の次元増大率$2\log\Lambda$を得た。
\texttt{verify\_boundary\_thresholds.py}はLaurent帰納指数、整数行列、
半群・固有値多重集合と有限範囲の数値を厳密に再現する。
全$n$の理論とMF定理を有限実験だけで証明したとはせず、
個別の適用の新規性も網羅的には確認していない。
\end{remark}

\section{境界の解析的modulus、Hochschild積、二捩れ付き楕円曲線}
\label{sec:boundary-analytic-moduli}
境界の二枝は、同じトーラス絡み目を持ちながら解析的には動き得る。
本節では、実際の反復曲線に残る係数不変量と、その重み付き主部に付随する
matrix factorization圏を区別して調べる。
基礎体は$\mathbb C$、$n\ge1$とし、Pell列の互いに素な整数
$2\le b_n<c_n$を用いる。$\kappa^2\ne1$、
$\rho_\pm=\kappa\pm\sqrt{\kappa^2-1}$と置く。

主論文の境界二枝定理から、条件を満たす初期conic $\{v=\kappa\}$の
反復逆像の境界germは、正則局所座標で
\begin{equation}\label{eq:boundary-weighted-model}
 f_{n,\kappa}=F_{n,\kappa}+\text{高重み項},\qquad
 F_{n,\kappa}=y^{2b_n}-2\kappa x^{c_n}y^{b_n}+x^{2c_n}
\end{equation}
と規格化できる。重みは$\operatorname{wt}(x)=b_n$、
$\operatorname{wt}(y)=c_n$である。
$F_{n,\kappa}$を\emph{重み付き模型}と呼び、一般の$f_{n,\kappa}$を
この多項式へ解析的に置き換えない。後半には、その置換が実際に不可能な例を示す。

\subsection{位相型が忘れる係数を解析同型から回収する}
\begin{proposition}[二枝の先頭係数の不変量]
\label{prop:boundary-kappa-analytic-invariant}
同じ互いに素な$2\le b<c$に対し、最低重み式が
$(y^b-r_1x^c)(y^b-r_2x^c)$となる二枝germを考える。
$r_1,r_2$は非零かつ相異なるとする。
二つのこの種のgermが解析的に同型なら
\[
 \mathcal I=\frac{(r_1+r_2)^2}{4r_1r_2}
\]
は一致する。特に実際の境界germの間の解析同型は
$\kappa^2=(\kappa')^2$を必要とする。
重み付き模型同士では、この条件は十分でもある。
\end{proposition}
\begin{proof}
唯一の接線は$y=0$なので、解析座標変換を$(X,Y)$とすると
$X=\alpha x+\text{高重み項}$、$\alpha\ne0$となる。
各枝で$\operatorname{ord}x=b$、$\operatorname{ord}y=c$である。
$Y$に$ib<c$の純粋項$x^i$があれば、それらの最小のものは枝上で
他の項と相殺せず、像の$Y$の位数が$c$未満になる。
従ってこれらの項は存在しない。
$b\nmid c$より重み$c$の単項式は$y$のみなので、
$Y=\beta y+\text{高重み項}$、$\beta\ne0$を得る。
各枝の値$y^b/x^c$は共通の定数$\beta^b/\alpha^c$倍されるだけである。
二枝の交換も許した上で、表示した量は不変となる。

規格化$r_1r_2=1$、$r_1+r_2=2\kappa$では$\mathcal I=\kappa^2$。
逆に$\kappa'=\kappa$なら同じ式であり、$\kappa'=-\kappa$なら
$\zeta^c=-1$を満たす$\zeta$を取り、$x\mapsto\zeta x$で
$F_{n,\kappa}(\zeta x,y)=F_{n,-\kappa}(x,y)$となる。
\end{proof}

重み付き模型の解析分類は、quasi-homogeneous平面曲線の一般分類の
二枝の場合でもある。\footnote{
L. M. C\^amara, \emph{On the classification of quasi-homogeneous curves},
\href{https://arxiv.org/abs/1009.1664}{arXiv:1009.1664v1}, Theorem 1(iii), \S4.3。
同論文の二点配置の$\mathbb C^*$商を、本稿の規格化に移したものとして比較する。}
本稿では実際の高重み項を持つgermにも適用できる\emph{必要条件}を直接証明した。
高重み項がある場合に$\kappa^2$だけで解析型が決まるとは主張しない。

\subsection{同じHochschild次元でも、積が圏を区別する}
\begin{theorem}[重み付き模型のdg Morita分類]
\label{thm:weighted-boundary-Morita-classification}
\[
 \mathcal Q_{n,\kappa}
 =\operatorname{MF}(\mathbb C[[x,y]],F_{n,\kappa})
\]
を二周期dg圏とする。必要ならそのperfect hullを取る。
このとき
\[
 \boxed{\mathcal Q_{n,\kappa}\simeq_{\rm Morita}\mathcal Q_{m,\kappa'}
 \quad\Longleftrightarrow\quad n=m\ \text{かつ}\ \kappa^2=(\kappa')^2.}
\]
これは\emph{重み付き模型}の分類であって、一般の実際の境界germ
$f_{n,\kappa}$の二周期MF圏の分類ではない。
\end{theorem}
\begin{proof}
$F_{n,\kappa}$の二偏微分は原点以外に共通零点を持たない。
Dyckerhoffの公式\footnote{
T. Dyckerhoff, \emph{Compact generators in categories of matrix factorizations},
Duke Math. J. 159 (2011), 223--274,
\href{https://doi.org/10.1215/00127094-1415869}{doi:10.1215/00127094-1415869};
\href{https://arxiv.org/abs/0904.4713v5}{arXiv:0904.4713v5},
Corollary 6.5（cohomologyを\emph{代数}として同定）、Theorem 5.9。
Theorem 6.6のhomologyの次元だけを用いるのではない。}
によって、Morita同値はJacobian環の$\mathbb C$-代数同型を与える。
その次元は$(2b_n-1)(2c_n-1)$で、Pell列に沿って真に増加するから$n=m$である。

重み付きEuler恒等式
\[
 b_nx(F_{n,\kappa})_x+c_ny(F_{n,\kappa})_y
       =2b_nc_nF_{n,\kappa}
\]
により、ここではJacobian環とTjurina環が等しい。
Mather--Yauの定理\footnote{
J. N. Mather, S. S.-T. Yau,
\emph{Classification of isolated hypersurface singularities by their moduli algebras},
Invent. Math. 69 (1982), 243--251,
\href{https://doi.org/10.1007/BF01399504}{doi:10.1007/BF01399504}, \S1の定理。
有限次元Tjurina環は解析環でも形式完備化でも同じである。}
は同じ埋込み次元の孤立複素超曲面germの接触同値を与える。
前命題から$\kappa^2=(\kappa')^2$が必要となる。
逆は前命題の明示的な\emph{右同値}による変数変換から従う。
\end{proof}

\begin{corollary}[固定された位相型の内部の連続な圏的modulus]
$n$を固定すると、全$\kappa^2\ne1$でlinkは$T(2b_n,2c_n)$、
Milnor数・Tjurina数は$(2b_n-1)(2c_n-1)$、
対数標準閾値は$(b_n+c_n)/(2b_nc_n)$である。
これらが全て同じでも、$\kappa^2$が異なれば上の二周期dg圏はMorita非同値である。
\end{corollary}

さらに生成元そのもののcohomology代数だけでも、この違いは消える。
安定化剰余体$G_{n,\kappa}$のendomorphism dg代数のminimal $A_\infty$模型は、
二次項が零であるため、基礎の積$m_2$として
$\bigwedge_{\mathbb C}(\theta_x,\theta_y)$を持つ。
しかしDyckerhoff Theorem 5.9の適合した模型では、符号規約を除いて
\[
 m_{b_n+c_n}(\underbrace{\theta_x,\ldots,\theta_x}_{c_n},
                     \underbrace{\theta_y,\ldots,\theta_y}_{b_n})
       =\pm(-2\kappa)
\]
に係数が現れる。これは選んだ模型での式であって、係数自体を
$A_\infty$座標非依存の不変量とする主張ではない。
Morita分類は上のHochschild\emph{積}とMather--Yauから得た。
生成元の通常のendomorphism cohomology代数をdg代数全体と同一視してはならない。

\subsection{四点配置からの楕円曲線とlevel $2$}
重み$(b_n,c_n)$の例外曲線のtorus座標$z=y^{b_n}/x^{c_n}$上では、
二枝の交点は$\rho_+,\rho_-$である。
さらにtoric境界の二端$0,\infty$を含めて四点を取り、
その二重被覆を
\begin{equation}\label{eq:boundary-level-two-elliptic}
 E^{\rm br}_\kappa:\quad \eta^2=z(z^2-2\kappa z+1)
\end{equation}
と定める。無限遠を原点$O$とし、非零二捩れ点$T=(0,0)$を標識する。
これは\emph{解消上の四点配置に付随させた楕円曲線}であり、
元の三次曲面の反標準楕円境界、Milnor fiber、あるいはMF圏の部分圏と
同一視するものではない。

\begin{proposition}[解析modulusは二捩れ付き楕円modulusとなる]
\label{prop:boundary-elliptic-level-two}
$E^{\rm br}_\kappa$は$\kappa^2\ne1$で滑らかであり、
$(E^{\rm br}_\kappa,O,T)$の同型類は$K=\kappa^2$で完全に決まる。
逆に原点と非零二捩れ点を標識した任意の複素楕円曲線がこの形で得られる。
従って$K\in\mathbb A^1\setminus\{1\}$は、非尖点の粗moduli曲線$Y_0(2)$の座標である。
その$j$不変量と判別式は
\begin{equation}\label{eq:boundary-elliptic-j}
 \Delta=64(K-1),\qquad
 \boxed{j=\frac{64(4K-3)^3}{K-1}}.
\end{equation}
$J_2=64(K-1)$と書けば$j=(J_2+16)^3/J_2$となる。
\end{proposition}
\begin{proof}
三つの有限分岐点$0,\rho_+,\rho_-$は相異なり、
$O$を加えた二重被覆はgenus $1$である。
一般の標識付き曲線を$Y^2=X(X^2+aX+b)$、$b\ne0$と書き、
$X=s^2z,Y=s^3\eta$、$s^4=b$で規格化すれば表示式になる。
規格化後に原点と$T$を保存するWeierstrass同型は
$z\mapsto s^2z$、$\eta\mapsto s^3\eta$、$s^4=1$である。
従って$\kappa$は符号を除いて定まり、$K$が完全不変量となる。
$\Delta$と$j$は$a_2=-2\kappa,a_4=1$から
$c_4=16(4\kappa^2-3)$を計算して得る。
$Y_0(2)$はまさに非零二捩れ点（位数2の部分群）を持つ楕円曲線の粗moduliである。
\end{proof}

\begin{corollary}[$j$だけでは圏を区別できない]
$(b,c)=(2,5)$、$\kappa=0$と$\kappa'=3\sqrt2/4$では
$j(E^{\rm br}_\kappa)=j(E^{\rm br}_{\kappa'})=1728$。
両MF圏のHochschild次元も$27$であるが、$K=0$と$K'=9/8$は異なるため
重み付き模型の圏はMorita非同値である。
\end{corollary}
標識を忘れた楕円曲線の$j$と、二捩れ標識を残した$Y_0(2)$の情報量の差である。
逆に$n$を変えて$\kappa$を固定すると、この楕円曲線は変わらないが
MF圏のHochschild次元はPell型に増大する。

\begin{proposition}[二同種写像はMF同値を意味しない]
\label{prop:boundary-two-isogeny}
$T$による商は次の明示的な次数2のisogenyで与えられる：
\[
 (z,\eta)\longmapsto
 \left(X=z-2\kappa+z^{-1},\quad Y=\eta(1-z^{-2})\right),
\]
\[
 Y^2=X\bigl(X^2+4\kappa X+4(\kappa^2-1)\bigr).
\]
非零二捩れ点$(X,Y)=(0,0)$を再び標識して規格化すると、
\[
 K\longmapsto K^\vee=\frac K{K-1},\qquad
 J_2\longmapsto\frac{4096}{J_2}
\]
となる。例えば$K=4$と$K^\vee=4/3$の付随楕円曲線は二同種だが、
同じ$n$の重み付きMF模型はMorita非同値である。
\end{proposition}
\begin{proof}
$T$による平行移動は$(z,\eta)\mapsto(z^{-1},-\eta/z^2)$。
表示した$X,Y$は不変で、直接代入して目標方程式を満たす。
関数体の次数が2なので、この平行移動による商である。
$X^2$の係数の平方と$X$の係数から規格化パラメータを計算すると
$K^\vee=K/(K-1)$となる。これはinvolutionであり、$J_2$の式も従う。
最後は定理\ref{thm:weighted-boundary-Morita-classification}による。
\end{proof}

このisogenyは標準的な二捩れ商の具体式であり、新しい一般isogeny理論ではない。
境界の解析的modulus、Jacobian環の積、二捩れ付き楕円曲線が
同じ$K$を記録するという比較に意味がある。

\subsection{実際の境界germは主部と異なる：$\mu=27,\ \tau=24$の例}
\begin{proposition}[一反復ですでに生じる高重み情報]
\label{prop:actual-boundary-tjurina-defect}
$A=B=C=0,h=1$の滑らかなLG曲面で、初期conicを$v=2$とする。
一反復$C_1=g^{-1}\{v=2\}$の境界germは
\[
 \mu(C_1,p_w)=27,\qquad \tau(C_1,p_w)=24
\]
を持つ。従って$y^4-4x^5y^2+x^{10}$という重み付き主部へ
解析的に接触同値ではない。
\end{proposition}
\begin{proof}
この初期conicは$\kappa^2-1=3$、$N_v(\kappa)=18$で条件を満たす。
$p_w$の射影chartで$x=U/W,y=V/W,t=T/W$とすると曲面式は
\[
 P=2xy-t(x^2+y^2+1)-2t^3=0.
\]
$r_1=2uv-w,r_2=2vr_1-u,r_3=2vr_2-r_1$、
$V_g=2r_2r_3-v$と置く。これは$g^*v$の正確な式である。
$H=t^7(V_g(x/t,y/t,1/t)-2)$は多項式で、
\[
 F(x,y)=\frac{\operatorname{Res}_t(P,H)}{64x^2y^5}
\]
も多項式となる。$P_t(0)=-1$であり、$t=0$以外の二つの局所根では
$H(0,0,t)=-2t^7$が非零なので、その二factorは単元である。
$x^2y^5$は境界の固定成分であるから、$F$は求める境界germの
局所被約方程式を単元倍したものとなる。

重み$(2,5)$で、計算に必要な重み$31$以下のjetは
\begin{align*}
 F_{\le31}={}&4(x^{10}-4x^5y^2+y^4)
 -16x(x^{10}-7x^5y^2+y^4)\\
 &+8x^2(3x^{10}-32x^5y^2+13y^4)
 -16x^3(x^{10}-17x^5y^2+21y^4)\\
 &+4x^4(x^{10}-64x^5y^2+227y^4)
 +2y^2(287x^{10}-1204x^5y^2+59y^4).
\end{align*}
先頭部分の二偏微分は重み$18,15$の正則列である。
filtered Koszul複体のexactnessにより、局所Jacobian環の随伴次数環は
この先頭二偏微分の商であり、最大socle重みは$18+15-2-5=26$となる。
従って重み$26$を超える単項式を切り捨てて\emph{正確に}計算できる。
偏微分に必要なjetが重み$31=26+5$までである理由もこれである。

重み$26$以下の45単項式を用いて、$F_x,F_y$の単項式倍の係数行列を作ると
rankは$18$で、商の基底を
\[
 \{1,x,\ldots,x^{13}\}\ \cup\
 \{y,xy,\ldots,x^8y\}\ \cup\
 \{y^2,xy^2,x^2y^2,x^3y^2\}
\]
と取れる。従って$\mu=45-18=27$。
この基底での有理数上の行消去は
\[
 [F]=12x^{13}-\tfrac{96}{5}x^{12}-\tfrac{12}{5}x^{11},\quad
 [xF]=-\tfrac{96}{5}x^{13}-\tfrac{12}{5}x^{12},\quad
 [x^2F]=-\tfrac{12}{5}x^{13},
\]
および$[x^3F]=[yF]=0$を与える。
従って$([F])$の次元は正確に3、$\tau=27-3=24$である。
表示したresultant、有限jet、行列rank、正規形は
\texttt{verify\_boundary\_analytic\_moduli.py}が元のVieta式から厳密に再生成する。
主部のTjurina数はEuler恒等式により$27$なので、接触同値は不可能である。
\end{proof}

\begin{remark}[比較の限界と次の問題]
この例では$\mu$とlinkは主部と同じでも、Tjurina環は違う。
一般の高重み項を持つ$F$について、二周期MFのHochschild環は
Jacobian環であり、Tjurina環はその中の\emph{指定した類}$[F]$で割った商である。
上の$\tau$の差だけから、標識なし二周期dg圏の非同値を結論してはいない。
実際の境界germへ重み付き模型の完全分類を拡張するには、
$[F]$に対応する追加構造、または別の圏的不変量を検討する必要がある。

一方、$\kappa^2$の解析的不変性は実際の境界にも適用される。
固定した$h=1,A=B=C=0$と$n$において、許容する$\kappa$を動かすだけで
同じアフィン曲線$\mathbb G_m$・同じlink型・同じ閾値を持つ無限個の
解析的に異なる境界が得られる。
本節の新規性の網羅確認は未了であり、quasi-homogeneous分類、
Mather--Yau、Dyckerhoff、楕円曲線の標準理論自体は既知として扱う。
\end{remark}

\section{境界の楕円modulusをMilnor被覆のPrym部分へ実現する}
\label{sec:boundary-milnor-prym}
前節の楕円曲線は、解消上の四点配置から付随させたものであった。
本節では、その曲線へ実際の有限写像を構成する。
ただし対象は\emph{重み付き主部のMilnor fiberに二次Kummer被覆を施したもの}である。
元のMilnor fiberから同じ底座標を保って写すことには、明示的な平方類の障害がある。
最後に、実際の高重みgermの復元について、整数次数付きdg特異点圏で得られる帰結を区別して述べる。

以下$b=b_n,c=c_n$、$n\ge1$とし、$b$は偶数、$c$は奇数、$\gcd(b,c)=1$を用いる。
$\kappa^2\ne1$、$K=\kappa^2$と置き、
\[
 V_{b,c;\kappa}=\{y^{2b}-2\kappa x^cy^b+x^{2c}=1\},\qquad
 E_\kappa:\ \eta^2=z(z^2-2\kappa z+1)
\]
とする。$C_{b,c;\kappa}$は$V_{b,c;\kappa}$の滑らかな射影模型、
$E_\kappa$では無限遠を原点とする。
重み付き斉次性により$V$は局所Milnor fiberの代数的模型である。
高重み項を持つ実境界の解析的Milnor fiberと無条件に同一視しない。

\subsection{元のfiberでの障害と、最小の二次被覆}
\begin{theorem}[平方類を解消するKummer被覆]
\label{thm:boundary-kummer-prym}
$C=C_{b,c;\kappa}$上の有理関数$z=y^b/x^c$を固定する。
$C\to E_\kappa$でこの$z$を引き戻す有理写像は存在しない。
一方、$\widetilde C$を$\mathbb C(C)(a)$、$a^2=x$の滑らかな射影模型とすると、
\begin{equation}\label{eq:milnor-elliptic-map}
 q:\widetilde C\longrightarrow E_\kappa,\qquad
 z=\frac{y^b}{a^{2c}},\qquad \eta=\frac{y^{b/2}}{a^{3c}}
\end{equation}
は次数$2bc$の有限写像となる。
$\pi:\widetilde C\to C$は連結な二重被覆で、分岐点は$x=0$の$2b$点である。
従って
\[
 g(C)=2bc-b-c,\qquad
 g(\widetilde C)=4bc-b-2c-1,\qquad
 \dim\operatorname{Prym}(\widetilde C/C)=2bc-c-1.
\]
\end{theorem}
\begin{proof}
$Q(z)=z^2-2\kappa z+1$とすると、$V$の方程式は$x^{2c}Q(z)=1$である。
$bp+cq_0=1$を満たす整数$p,q_0$を選び、$t=x^py^{q_0}$と置くと
\[
 x=t^bz^{-q_0},\quad y=t^cz^p,\qquad
 t^{2bc}=\frac{z^{2cq_0}}{Q(z)}.
\]
右辺は$Q$の二根で位数$-1$を持つので、このKummer拡大は連結で次数$2bc$。
四つの底点$0,\infty,\rho_+,\rho_-$上の点数は、それぞれ$2c,2b,1,1$である。
Riemann--Hurwitzから$g(C)=2bc-b-c$を得る。
$\rho_\pm$上の二点では$x$の極位数が$b$であり、$z=\infty$上の$2b$点では
$x$の零位数が$bp+cq_0=1$である。
これは$x=0,y^{2b}=1$という$V$の$2b$点でも直接確認できる。

$E_\kappa$へ持ち上げるための関数の平方は
\begin{equation}\label{eq:milnor-square-obstruction}
 zQ(z)=\frac{y^b}{x^{3c}}
      =x\left(\frac{y^{b/2}}{x^{(3c+1)/2}}\right)^2
\end{equation}
である。$x$は単純零点を持つため平方ではない。これが元の$C$での障害である。
同式により平方根を添加する拡大はちょうど$\mathbb C(C)(\sqrt x)$となり、
表示\eqref{eq:milnor-elliptic-map}を得る。滑らかな射影曲線間の有理写像は延長される。
$\widetilde C\to\mathbb P^1_z$の次数は$4bc$、$E_\kappa\to\mathbb P^1_z$は次数2なので$q$の次数は$2bc$。
$x$の極位数$b$は偶数であるから、$\pi$は無限遠では分岐せず、$2b$個の単純零点だけで分岐する。
再びRiemann--Hurwitzを使い、残る数値を得る。
\end{proof}

\begin{corollary}[楕円曲線のPrym因子]
\label{cor:elliptic-prym-factor}
$E_\kappa$は$\operatorname{Prym}(\widetilde C/C)$のisogeny因子である。
より正確には$q^*$の像はPrymに含まれ、
\[
 e_\kappa=\frac{1}{2bc}q^*q_*
 \in\operatorname{End}^0(J(\widetilde C))
\]
はその像への冪等射である。
従って$H^1(E_\kappa,\mathbb Q)$は$H^1(\widetilde C,\mathbb Q)$の
deck involutionに関する反不変部分のHodge直和因子となる。
\end{corollary}
\begin{proof}
$\delta(a,y)=(-a,y)$は$q$の座標では$(z,\eta)\mapsto(z,-\eta)$となる。
後者は$E_\kappa$の逆元写像なので$\delta^*q^*=-q^*$。
Prymは$\ker(1+\delta^*)^0$であり、連結な$q^*(E_\kappa)$はそこへ入る。
曲線のnormとpullbackについて$q_*q^*=[2bc]$であるから$q^*$は有限核を持ち、
表示した$e_\kappa$は冪等である。
これはPrymの標準的なnorm理論\footnote{
D. Mumford, \emph{Prym Varieties I}, in \emph{Contributions to Analysis},
Academic Press (1974), 325--350, \S\S1--3。
normとpullback、および$P=\ker(1+\delta)^0$については特に印刷頁331を参照。
}のこの被覆への適用であり、整数係数や主偏極を保つ積分解を主張するものではない。
\end{proof}

\subsection{全てのPell指数に共通するCiani四次曲線}
\begin{theorem}[genus三の中間模型と三楕円分解]
\label{thm:universal-ciani-milnor}
$\widetilde C$から、$n$に依存しない滑らかな平面四次曲線
\[
 \mathcal C_\kappa:\quad X^4+Y^4-2\kappa X^2Y^2=Z^4
\]
への次数$bc/2$の写像が、アフィンchartで
\[
 (a,y)\longmapsto(X,Y)=(a^c,y^{b/2})
\]
により与えられる。
$\alpha(X,Y)=(-X,Y)$、$\beta(X,Y)=(X,-Y)$とすると三つの商は楕円曲線であり、
\[
 J(\mathcal C_\kappa)\sim E_i^2\times D_\kappa,\qquad
 D_\kappa:\ T^2=r^4-2\kappa r^2+1,
\]
となる。ここで$E_i$は$j=1728$の楕円曲線、$\sim$は複素数体上のisogenyである。
さらに$\operatorname{Prym}(\widetilde C/C)$は$E_i\times D_\kappa$にisogenousな
二次元の部分アーベル多様体を持つ。
\end{theorem}
\begin{proof}
式への代入で写像を得る。一般の$X,Y\ne0$には、$a^c=X$の$c$通りと
$y^{b/2}=Y$の$b/2$通りの持上げがあり、次数は$bc/2$。
$\kappa^2\ne1$で平面四次は滑らかなので、そのgenusは3である。

$\alpha$商と$\beta$商はそれぞれ
\[
 (X^2-\kappa Y^2)^2=1+(K-1)Y^4,\qquad
 (Y^2-\kappa X^2)^2=1+(K-1)X^4
\]
であり、いずれも複素数体上で$j=1728$となる。
$\alpha\beta$商は$r=Y/X,T=1/X^2$による$D_\kappa$である。
三つの商写像による$H^0(\Omega^1)$の引戻しは、Klein四元群の相異なる三つの非自明指標に属し、
各次元1で全体を張る。従って三楕円曲線からのJacobianのpullbackの和はisogenyとなる。
これはCiani四次の標準的な三楕円分解の特殊例である。\footnote{
E. W. Howe, F. Lepr\'evost, B. Poonen,
\emph{Large torsion subgroups of split Jacobians of curves of genus two or three},
Forum Math. 12 (2000), 315--364,
\href{https://arxiv.org/abs/math/9809210}{arXiv:math/9809210}, \S4.1。
本稿では三つの商の方程式と文字分解を直接確認する。
}

$\widetilde C\to\mathcal C_\kappa$は$\delta$と$\alpha$に関して等変である。
$\alpha$反不変部分は$\beta$商と$\alpha\beta$商の二楕円因子であり、
そのpullbackは有限核を持って$\operatorname{Prym}(\widetilde C/C)$へ入る。
これが最後の主張である。
\end{proof}

$D_\kappa$から前節の$E_\kappa$へは
\begin{equation}\label{eq:ciani-level2-isogeny}
 (r,T)\longmapsto(z,\eta)=(r^2,rT)
\end{equation}
という次数2の写像がある。
そのdeck変換$(r,T)\mapsto(-r,-T)$は固定点を持たず、
原点を適合させれば二同種写像となる。
\[
 j(D_\kappa)=\frac{64(K+3)^3}{(K-1)^2}
 =j\left(E_{\kappa^\vee}\right),\qquad
 (\kappa^\vee)^2=\frac K{K-1}.
\]
従って前節のlevel $2$ isogenyは、ここでは\emph{共通のgenus三模型の商写像}として実現される。
$n=1$では$g(C)=13$、$g(\widetilde C)=27$、Prym次元14、
$\widetilde C\to\mathcal C_\kappa$の次数5、$\widetilde C\to E_\kappa$の次数20である。
高genusの増大と独立に、可変な楕円因子は同じ式で残る。

\subsection{Prym部分の周期はLegendre型方程式を満たす}
\begin{proposition}[可変楕円因子のPicard--Fuchs方程式]
\label{prop:boundary-prym-PF}
$E_\kappa$上の形式$\omega_\kappa=dz/\eta$と平坦に動くcycleの周期$\Pi$は
\begin{equation}\label{eq:boundary-prym-PF}
 (\kappa^2-1)\Pi''+2\kappa\Pi'+\tfrac14\Pi=0
\end{equation}
を満たす。同式は系\ref{cor:elliptic-prym-factor}で得たPrym部分のrank二の周期系にも適用される。
$q=(1-\kappa)/2$とすれば、これは
$q(1-q)\Pi_{qq}+(1-2q)\Pi_q-\Pi/4=0$、すなわち
パラメータ$(1/2,1/2;1)$のGauss超幾何方程式である。
\end{proposition}
\begin{proof}
$P=z(z^2-2\kappa z+1)$と書く。
$\partial_\kappa\omega=z^2dz/P^{3/2}$、
$\partial_\kappa^2\omega=3z^4dz/P^{5/2}$なので、直接計算で
\[
 \left((\kappa^2-1)\partial_\kappa^2+2\kappa\partial_\kappa+\tfrac14\right)
 \omega_\kappa
 =d_z\left(\frac{z^2(1-z^2)}{2P^{3/2}}\right)
\]
を得る。右辺は楕円曲線の有理関数の微分であり、閉cycleでの積分は零となる。
有限写像によるpullbackとtraceはGauss--Manin接続に適合するため、Prymの直和因子にも同じ式を得る。
\end{proof}

これは$\kappa$という\emph{解析modulusを動かす}周期系である。
固定$\kappa$で値$F$を一周する有限位数のMilnor monodromyや、
元のLG三次曲面のlevel $h$を動かしたarcsine周期系とは異なる。
前節で$j$やisogenyが二周期MF模型を完全には分類しなかったこととも矛盾しない。
ここで回収したのは、主部のMilnor被覆の\emph{有理Hodge因子}であって、圏の同値ではない。

\subsection{実germの高重み情報を整数次数付きdg圏で検出する}
最後に実際の$\mu=27,\tau=24$例へ戻る。
$R_f=\mathbb C[[x,y]]/(f)$とし、
\[
 \mathcal S_f=D^b_{\rm dg}(\operatorname{mod}R_f)/\operatorname{Perf}_{\rm dg}(R_f)
\]
を、その通常の$\mathbb Z$-graded dg enhancementを持つ特異点圏とする。
本節のHochschild cohomologyは、この整数次数付き圏の絶対的なものとする。

\begin{corollary}[実境界と主部のdg特異点圏は非同値]
\label{cor:actual-germ-Zdg-distinction}
前節の一反復germ $f$と、その重み付き主部$F$について
\[
 \dim\mathrm{HH}^0(\mathcal S_f)=24,\qquad
 \dim\mathrm{HH}^0(\mathcal S_F)=27.
\]
従って$\mathcal S_f$と$\mathcal S_F$はdg Morita非同値である。
さらに、同じ$n$の許容する実境界germ同士で$\kappa^2\ne(\kappa')^2$なら、
それらの整数次数付きdg特異点圏もMorita非同値である。
\end{corollary}
\begin{proof}
Hua--Kellerの復元定理\footnote{
Z. Hua, B. Keller, \emph{Cluster categories and rational curves},
Geom. Topol. 28 (2024), 2569--2634,
\href{https://doi.org/10.2140/gt.2024.28.2569}{doi:10.2140/gt.2024.28.2569};
\href{https://arxiv.org/abs/1810.00749v6}{arXiv:1810.00749v6}, Theorem 5.9, pp.36--37。
ここで使うのは一般の孤立超曲面についての定理であり、
三次元flopやNCCRの存在を仮定する後続の定理ではない。
}は、孤立超曲面germについて
$\mathrm{HH}^0(\mathcal S_f)\cong\mathbb C[[x,y]]/(f,f_x,f_y)$を代数として与える。
前節の厳密なTjurina数の計算とMorita不変性から最初の主張が従う。
後半では、Morita同値ならTjurina環が同型となり、Mather--Yauによって接触型が一致する。
前節で証明した実二枝の解析的不変量$\kappa^2$が一致しなければならない。
\end{proof}

\begin{remark}[基礎環・次数・実germを保つ]
二周期MFの相対的Hochschild環をJacobian環と同定する公式と、
上の整数次数付きdg特異点圏の絶対的$\mathrm{HH}^0$をTjurina環と同定する公式を入れ替えない。
本系はHua--Kellerの指定するenhancementでの結論であり、
本稿ではperiodization・忘却の全比較を構成せず、一般の二周期MFの復元へ無条件に言い換えない。
同様に、Kummer--Prym構成は重み付き模型のMilnor curveについて証明したもので、
$\tau=24$の実際のgermに解析的に移したものではない。

Prymのnorm理論、Cianiの三楕円分解、超幾何方程式、dg復元の一般定理は既知である。
本稿の位置付けは、境界から得た同じ$K$を明示的な有限写像とHodge因子へ結び、
何を付加すれば幾何学的対応が実現し、どの範囲にはまだ移せないかを示すことにある。
Prym全体の偏極を保つ直積分解、実germのPrym比較、HMSは未構成である。
\end{remark}

\section{巡回Milnor被覆の整数偏極と極端楕円曲面}
\label{sec:prym-integral-polarization}
前節で得た$C,\widetilde C,D_\kappa,E_\kappa$と写像を、
有限Galois被覆および整数偏極の観点から精密化する。
$b=b_n=2m,\ c=c_n$はPell列の互いに素な正整数とし、
$n\ge1$、$c$は奇数、$\kappa^2\ne1$とする。
$K=\kappa^2$と書く。前節の方程式を再掲すると
\[
 \widetilde C:\quad y^{4m}-2\kappa a^{2c}y^{2m}+a^{4c}=1,
 \qquad
 D_\kappa:\quad T^2=r^4-2\kappa r^2+1,
\]
\[
 E_\kappa:\quad \eta^2=z(z^2-2\kappa z+1)
\]
である。曲線はすべて滑らかな射影模型を意味する。
\begin{equation}\label{eq:cyclic-prym-factorization}
 \widetilde C\xrightarrow{\,q'\,}D_\kappa
 \xrightarrow{\,\psi\,}E_\kappa,\qquad
 q'(a,y)=\left(\frac{y^m}{a^c},\,\frac1{a^{2c}}\right),
 \quad \psi(r,T)=(r^2,rT)
\end{equation}
とおく。合成は前節の$ q:\widetilde C\to E_\kappa$である。

\begin{theorem}[巡回分岐と最大非分岐商]
\label{thm:prym-cyclic-inertia}
$q'$は次数$bc$の巡回Galois被覆、$q$は次数$2bc$の
巡回Galois被覆である。
$q'$の分岐は正確に次の三種類で記述される。
\[
\begin{array}{c|c|c|c}
 D_\kappa\text{上の点}&\text{底の点数}&
 q'\text{の各fiberの点数}&\text{分岐指数}\\ \hline
 T=0&4&1&bc\\
 r=0&2&2c&m=b/2\\
 r=\infty&2&b&c
\end{array}
\]
各行の分岐指数が$1$なら、その行は実際には分岐しない。
とくに$T=0$の四点は完全分岐である。
$q$の最大非分岐Galois中間被覆は
$D_\kappa\to E_\kappa$であり、その次数は$2$である。
\end{theorem}
\begin{proof}
$\gcd(m,c)=1$である。
$\zeta_{2c},\zeta_{2m}$を原始根とすると
\[
 (a,y)\longmapsto(\zeta_{2c}a,\zeta_{2m}y)
\]
は$\widetilde C$の方程式を保ち、$r=y^m/a^c$と$T=a^{-2c}$を固定する。
位数は$\mathrm{lcm}(2c,2m)=2cm=bc$で、既に証明した$q'$の次数と一致する。

さらに$\zeta$を原始$4c$乗根とし、
$j\equiv3\pmod4$かつ$\gcd(j,4m)=1$となる整数$j$を選ぶ。
$\xi=\exp(2\pi i j/(4m))$とすると
\[
 (a,y)\longmapsto(\zeta a,\xi y),\qquad
 \xi^m=\zeta^{3c},
\]
は$z=y^{2m}/a^{2c}$と$\eta=y^m/a^{3c}$を固定する。
その位数は$\mathrm{lcm}(4c,4m)=4cm=2bc$で、
$q$も巡回Galois被覆となる。

$q'$の分岐点は方程式から直接求まる。
$r=0,T=\pm1$では$y=0$、$a^{2c}=1/T$で、
各点の逆像は$2c$点、局所的に$r\sim y^m$なので指数$m$。
$r=\infty$の二点では$a=0$、各点に$y^{2m}=\pm1$の
$b=2m$個の逆像があり、指数$c$。
$T=0$の四点では$a,y$が極を持つ。
前節の$z=y^b/x^c$のKummer表示によって
各点の逆像はただ一つ、局所次数は$bc$となる。
これで表の全分岐が得られる。整合性の検算は
\[
 4(bc-1)+2(bc-2c)+2(bc-b)
   =2g(\widetilde C)-2
\]
というRiemann--Hurwitz式である。

$q$は巡回次数$2bc$であり、その中間の
$q'$は$T=0$上で完全分岐である。
したがって$q$の慣性部分群の生成する部分群は
ちょうど位数$bc$で、最大非分岐商はその指数$2$の商
$D_\kappa\to E_\kappa$である。
\end{proof}

\begin{theorem}[楕円因子の核と誘導偏極]
\label{thm:prym-elliptic-theta-degree}
$q^*:J(E_\kappa)\to J(\widetilde C)$の核は正確に
\[
 \boxed{\ker q^*
 =\bigl\{0,\,[\,(0,0)-O\,]\bigr\}
 \cong\mathbb Z/2\mathbb Z.}
\]
$q'^*:J(D_\kappa)\to J(\widetilde C)$は単射で、その像は
$\operatorname{Prym}(\widetilde C/C)$内の楕円部分多様体である。
$J(\widetilde C)$の標準theta偏極をこの楕円部分多様体へ
制限した偏極の次数は、正確に
\begin{equation}\label{eq:prym-exact-theta-degree}
 \boxed{\deg\bigl(\Theta_{J(\widetilde C)}|_{q'^*J(D_\kappa)}\bigr)=bc.}
\end{equation}
また、$C$のgenusを$g=2bc-b-c$とすれば、
$\operatorname{Prym}(\widetilde C/C)$へのtheta制限の
偏極型は$(1^{b-1},2^g)$である。
\end{theorem}
\begin{proof}
一般に、有限被覆$f:X\to Y$が完全分岐するfiberを持てば、
$f^*:J(Y)\to J(X)$は単射である。
実際、核の元$L$は$f_*f^*=[\deg f]$からtorsionとなる。
$L$が非自明なら、その有限位数に対応する連結な非自明
非分岐巡回被覆$Y'\to Y$を$f$が経由する。
しかし完全分岐fiberは一点しか持たず、
非分岐被覆の各fiberは複数点なので矛盾する。
定理\ref{thm:prym-cyclic-inertia}を$q'$に適用して単射を得る。

$E_\kappa$上の関数$z$のdivisorは
$2(0,0)-2O$なので、$D_\kappa$は
$\sqrt z$を添加した非自明な非分岐二重被覆である。
従って$\psi^*$の核は標識二捩れ類$[(0,0)-O]$が生成する群である。
$q^*=q'^*\psi^*$と$q'^*$の単射性から、最初の式が従う。

$\psi$は$D_\kappa$の非零二捩れに対応する平行移動の商なので、
$\psi^*$の像は$J(D_\kappa)$全体である。
前節で示したdeck involutionの反変性から、
$q'^*J(D_\kappa)=q^*J(E_\kappa)$はPrymに入る。
曲線間の次数$bc$の写像$q'$について
$q'_*q'^*=[bc]$が成立し、Jacobianの主偏極を用いて
$q'^*$のtheta引戻しは$bc$倍の楕円主偏極になる。
$q'^*$は単射なので、その像へのtheta制限の次数が$bc$となる。
なお$E_\kappa$上で直接引き戻すthetaの次数は$2bc$であり、
同じ数と混同しない。

最後に$\widetilde C\to C$は$2b$点で分岐する二重被覆である。
Mumfordのramified Prym偏極の標準公式
（前節の文献参照）を使うと、偏極型は$(1^{b-1},2^g)$となる。
これは楕円因子上の次数$bc$と区別すべき、Prym全体の偏極型である。
\end{proof}

\begin{corollary}[Pell力学がPrymの整数偏極に現れる]
\label{cor:prym-pell-theta-growth}
最初の三反復では
\[
 (b,c)=(2,5),(12,29),(70,169)
 \quad\Longrightarrow\quad
 \deg(\Theta|_{D_\kappa})=10,348,11830
\]
である。$\Lambda=3+2\sqrt2$とすれば
\[
 \lim_{n\to\infty}\frac1n
   \log\deg(\Theta|_{D_\kappa})=2\log\Lambda.
\]
同じ$E_\kappa$がPrymの有理Hodge因子として繰り返し現れても、
それが埋め込まれる際の整数偏極はPell指数とともに増大する。
\end{corollary}
\begin{proof}
式\eqref{eq:prym-exact-theta-degree}と
$b_n,c_n$のPell漸近式から直ちに従う。
\end{proof}

\subsection{Prymの楕円族は極端有理楕円曲面を与える}
\begin{proposition}[Miranda--Perssonの$X_{411}$とBelyi写像]
\label{prop:prym-extremal-X411}
$\kappa$を変数とする楕円族
\[
 \mathscr E:\quad \eta^2=z^3-2\kappa z^2+z
\]
の$\mathbb P^1_\kappa$上の相対的極小正則模型は、
特異fiber型$I_1,I_1,I_4^*$を持つ極端有理楕円曲面である。
Miranda--Perssonの記法では$X_{411}$に一致し、
Mordell--Weil群は$\mathbb Z/2\mathbb Z$で、その生成元は
$T=(0,0)$である。
$j$写像
\begin{equation}\label{eq:prym-belyi-j}
 j(\kappa)=\frac{64(4\kappa^2-3)^3}{\kappa^2-1}
\end{equation}
は次数$6$のBelyi写像であり、
$j=0,1728,\infty$上の分岐型はそれぞれ
\[
 (3,3),\qquad(2,2,2),\qquad(4,1,1)
\]
である。
\end{proposition}
\begin{proof}
Weierstrass不変量を直接計算すると
\[
 \Delta=64(\kappa^2-1),\qquad
 c_4=16(4\kappa^2-3),\qquad j=c_4^3/\Delta.
\]
$\kappa=\pm1$では$\Delta$が単純零点、$c_4\ne0$なので
両fiberは$I_1$である。
無限遠では$s=1/\kappa$、$z=s^{-2}Z$、$\eta=s^{-3}Y$とすると
\[
 Y^2=Z^3-2sZ^2+s^4Z,\qquad
 \Delta_\infty=64s^{10}(1-s^2),\quad
 c_{4,\infty}=16s^2(4-3s^2).
\]
特性零のKodaira--Tate表からこれは$I_4^*$である。
この方程式は基本束$\mathcal O_{\mathbb P^1}(1)$を持つ
最小Weierstrass模型なので、零sectionを持つ有理楕円曲面となる。
$I_4^*$のroot latticeは$D_8$でrank $8$だから、
Shioda--Tateと$\rho=10$によりMordell--Weil rankは$0$。
そのtorsionはMiranda--Perssonの分類と、
$T=(0,0)$の存在から正確に$\mathbb Z/2\mathbb Z$である。

最後に
\[
 j-1728=
 \frac{64\kappa^2(8\kappa^2-9)^2}{\kappa^2-1}
\]
である。$j=0$では$\kappa=\pm\sqrt3/2$が各三重、
$j=1728$では$\kappa=0,\pm\sqrt{9/8}$が各二重、
$j=\infty$では$\kappa=\infty$が四重、
$\kappa=\pm1$が各単純極となる。
従って分岐は三値だけで、表示した分岐型を得る。

この楕円曲面の分類とBelyi型は既知である。
Miranda--Persson,
\emph{On extremal rational elliptic surfaces},
Math.\ Z.\ \textbf{193} (1986), 537--558,
\href{https://doi.org/10.1007/BF01160474}{doi:10.1007/BF01160474},
Table~5.2（$X_{411}$）を一次資料で確認した。
本稿での役割は、この既知曲面がPell境界からの
Prym因子の変動として現れることの明示的な説明である。
\end{proof}

\begin{remark}[同じレベル二構造、異なる三つの幾何]
\label{rem:prym-three-levels}
標識二捩れ$T=(0,0)$は、単に四点配置に選んだ補助標識ではなく、
定理\ref{thm:prym-elliptic-theta-degree}の自然な
$J(E_\kappa)\to J(\widetilde C)$の核そのものとして回復される。
同じパラメータ$K=\kappa^2$は、
重み付きMF模型の解析的modulus、
楕円曲線の$Y_0(2)$標識、
Prymの二同種核を同時に記録する。
ただしMFの同値を楕円isogenyから得たわけではない。
さらに$\kappa$を動かすモジュラー楕円曲面のmonodromy、
固定$\kappa$におけるMilnorのmonodromy、
元のLG等位面の$h$に沿う周期は、互いに別の作用である。
\end{remark}


\section{残る比較問題}
第一に、実germの高重み項がMilnor被覆のHodge構造とPrym因子へ及ぼす影響を調べる必要がある。
$\mu=27,\tau=24$の例は、重み付き主部へ解析的に置き換える方法の限界を示す。
第二に、二周期MFの相対Hochschild構造と整数次数付きdg圏の絶対Hochschild構造の比較では、
基礎環、周期化、指定された関数の類$[f]$を保持しなければならない。
第三に、極filtrationの歪み率をmirror側のaction filtrationへ移すには、
圏同値だけでなく、指定した$g$との同変性と一様なfiltration比較が必要である。
本文の有限計算は多項式恒等式・整数式・指定した局所代数を検算するもので、
これらの未構成の比較や新規性の網羅調査を代替しない。

\UsePublicationReferences{BoundaryCamara,BoundaryCasals,BoundaryDHKK,BoundaryDyckerhoff,BoundaryGoldmanNeumann,BoundaryHowald,BoundaryHuaKeller,BoundaryMatherYau,BoundaryMirandaPersson,BoundaryMumford,CantatLorayCayley,CompanionFricke,JangBoundary}
\begin{thebibliography}{99}
\small
\linespread{0.94}\selectfont
\setlength{\itemsep}{0pt}
\ifcsname pubref@DeDekenLowen\endcsname
\bibitem{DeDekenLowen} O. De Deken, W. Lowen,
\emph{Abelian and derived deformations in the presence of $\mathbb Z$-generating geometric helices},
J. Noncommut. Geom. 5 (2011), 477--505,
\href{https://doi.org/10.4171/JNCG/83}{doi:10.4171/JNCG/83}, Theorem 4.5, Proposition 4.6, pp. 495--496.
\fi
\ifcsname pubref@StacksPi1\endcsname
\bibitem{StacksPi1} The Stacks Project Authors,
\emph{The Stacks Project},
\href{https://stacks.math.columbia.edu/tag/0BTX}{Tag 0BTX}, Lemma 58.14.3
(arithmetic--geometric fundamental group exact sequence).
\fi
\ifcsname pubref@StacksPerfect\endcsname
\bibitem{StacksPerfect} The Stacks Project Authors,
\emph{The Stacks Project},
\href{https://stacks.math.columbia.edu/tag/0BDK}{Tag 0BDK}, Lemma 36.31.3.
\fi
\ifcsname pubref@EG\endcsname
\bibitem{EG} P. Etingof, V. Ginzburg,
\emph{Noncommutative del Pezzo surfaces and Calabi--Yau algebras},
J. Eur. Math. Soc. 12 (2010), 1371--1416.
\href{https://arxiv.org/abs/0709.3593v4}{arXiv:0709.3593v4}.
Theorem 3.3.2(i),(ii), Proposition 7.3.1, Appendix 9.1。
\fi
\ifcsname pubref@LBS\endcsname
\bibitem{LBS} L. Le Bruyn, S. P. Smith, M. Van den Bergh,
\emph{Central extensions of three dimensional Artin--Schelter regular algebras},
Math. Z. 222 (1996), 171--212, Theorem 2.6, Corollary 2.7。
\fi
\ifcsname pubref@BG\endcsname
\bibitem{BG} A. Braverman, D. Gaitsgory,
\emph{Poincar\'e--Birkhoff--Witt theorem for quadratic algebras of Koszul type},
J. Algebra 181 (1996), 315--328, Theorem 0.5。
\fi
\ifcsname pubref@IS\endcsname
\bibitem{IS} N. Iyudu, S. Shkarin,
\emph{Three dimensional Sklyanin algebras and Gr\"obner bases},
\href{https://arxiv.org/abs/1601.00564}{arXiv:1601.00564}, Theorem 1.1, Corollary 6.1。
\fi
\ifcsname pubref@NS\endcsname
\bibitem{NS} C. Negron, T. Schedler,
\emph{The Hochschild cohomology ring of a global quotient orbifold},
Adv. Math. 364 (2020), 106978.
\href{https://arxiv.org/abs/1809.08715}{arXiv:1809.08715}, Theorem 1.1, Appendix A。
\fi
\ifcsname pubref@FIOS\endcsname
\bibitem{FIOS} E. Faber, C. Ingalls, S. Okawa, M. Satriano,
\emph{On stacky surfaces and noncommutative surfaces},
\href{https://arxiv.org/abs/2206.13359v2}{arXiv:2206.13359v2}, 8 February 2024,
Theorem 1.5, Corollary 1.7。Brauer捩れ付きの比較として使用する。
\fi
\ifcsname pubref@FPSS\endcsname
\bibitem{FPSS} L. Fu, M. Porta, S. Scherotzke, N. Sibilla,
\emph{Hochschild--Kostant--Rosenberg isomorphism for derived Deligne--Mumford stacks},
\href{https://arxiv.org/abs/2509.00501v3}{arXiv:2509.00501v3}, 20 January 2026,
Proposition 6.17。一般root stackの恒等部分への帰着を比較し、$n=2$ に限る後半の式を $n=3$ へ流用しない。
\fi
\ifcsname pubref@TerwilligerDAHA\endcsname
\bibitem{TerwilligerDAHA} P. Terwilliger,
\emph{The universal Askey--Wilson algebra and DAHA of type $(C_1^\vee,C_1)$},
SIGMA 9 (2013), 047, \href{https://doi.org/10.3842/SIGMA.2013.047}{doi:10.3842/SIGMA.2013.047},
Definitions 2.1, 2.2, 2.6; Lemmas 2.3, 2.5, 2.8.
\fi
\ifcsname pubref@GoldmanTrace\endcsname
\bibitem{GoldmanTrace} W. M. Goldman,
\emph{Trace coordinates on Fricke spaces of some simple hyperbolic surfaces},
in \emph{Handbook of Teichm\"uller Theory}, Vol. II (2009), 611--684,
\href{https://doi.org/10.4171/055-1/16}{doi:10.4171/055-1/16},
\S5.2.1, equation (5.2.1.2).
\fi
\ifcsname pubref@GoldmanToledo\endcsname
\bibitem{GoldmanToledo} W. M. Goldman, D. Toledo,
\emph{Affine cubic surfaces and relative $\mathrm{SL}(2)$-character varieties of compact surfaces},
\href{https://arxiv.org/abs/1006.3838}{arXiv:1006.3838}.
Theorem 1: every smooth-at-infinity cubic with generic tritangent plane is a four-holed-sphere character variety.
\fi
\ifcsname pubref@SchockD4\endcsname
\bibitem{SchockD4} N. Schock,
\emph{The $W(E_6)$-invariant birational geometry of the moduli space of marked cubic surfaces},
Math. Nachr. 297 (2024), 2639--2666,
\href{https://doi.org/10.1002/mana.202300459}{doi:10.1002/mana.202300459}.
Section 2: tritangent triangle and its orthogonal $D_4$ root lattice.
\fi
\ifcsname pubref@IISDiscriminant\endcsname
\bibitem{IISDiscriminant} M. Inaba, K. Iwasaki, M.-H. Saito,
\emph{Dynamics of the sixth Painlev\'e equation},
S\'eminaires \& Congr\`es 14 (2006), 103--167,
\href{https://emis.de/ft/47232}{original text}, \S9.1, Definition 9.1.
Fricke判別式をトレースパラメータへ引き戻した場合の二乗因子について参照する。
本稿の$N_5^2M_5$との同一視は主張しない。
\fi
\ifcsname pubref@HavlicekPosta\endcsname
\bibitem{HavlicekPosta} M. Havl\'{\i}\v{c}ek, S. Po\v{s}ta,
\emph{Center of quantum algebra $U'_q(\mathrm{so}_3)$},
J. Math. Phys. 52 (2011), 043521,
\href{https://doi.org/10.1063/1.3579992}{doi:10.1063/1.3579992}.
三乗根を含むroot-of-unity中心の構造について参照し、規約の完全な同一視は追加検証事項。
\fi
\ifcsname pubref@CantatLorayCayley\endcsname
\bibitem{CantatLorayCayley} S. Cantat, F. Loray,
\emph{Dynamics on Character Varieties and Malgrange irreducibility of Painlev\'e VI equation},
Ann. Inst. Fourier 59 (2009), 2927--2978,
\href{https://doi.org/10.5802/aif.2512}{doi:10.5802/aif.2512}, \S1.5.
\fi
\ifcsname pubref@JangBoundary\endcsname
\bibitem{JangBoundary} S. u. Jang,
\emph{Tropical Dynamics of Markov Surfaces},
\href{https://arxiv.org/abs/2306.11357v2}{arXiv:2306.11357v2},
Lemma 4.8 and Theorem 4.10. 取得した2025年12月版の題名と番号を使用する。
\fi
\ifcsname pubref@MikoszWeberTriality\endcsname
\bibitem{MikoszWeberTriality} M. Mikosz, A. Weber,
\emph{Triality, characteristic classes, $D_4$ and $G_2$ singularities},
J. Homotopy Relat. Struct. 10 (2015), 995--1011,
\href{https://doi.org/10.1007/s40062-014-0094-3}{doi:10.1007/s40062-014-0094-3}, Proposition 3, \S5.
\fi
\ifcsname pubref@DixonElliptope\endcsname
\bibitem{DixonElliptope} K. Dixon,
\emph{The multi-moment maps of the nearly K\"ahler $S^3\times S^3$},
Geom. Dedicata 200 (2019), 351--362,
\href{https://doi.org/10.1007/s10711-018-0375-9}{doi:10.1007/s10711-018-0375-9}, Lemma 3.5.
\fi
\ifcsname pubref@HikamiCluster\endcsname
\bibitem{HikamiCluster} K. Hikami,
\emph{Note on character varieties and cluster algebras},
SIGMA 15 (2019), 003,
\href{https://doi.org/10.3842/SIGMA.2019.003}{doi:10.3842/SIGMA.2019.003}.
\fi
\ifcsname pubref@AWAvatars\endcsname
\bibitem{AWAvatars} N. Cramp\'e et al.,
\emph{The Askey--Wilson algebra and its avatars},
J. Phys. A 54 (2021), 063001,
\href{https://doi.org/10.1088/1751-8121/abd783}{doi:10.1088/1751-8121/abd783}.
\fi
\ifcsname pubref@CookeLacabanne\endcsname
\bibitem{CookeLacabanne} J. Cooke, N. Lacabanne,
\emph{Higher rank Askey--Wilson algebras as skein algebras},
Ann. Inst. Fourier 76 (2026), 1055--1117,
\href{https://doi.org/10.5802/aif.3729}{doi:10.5802/aif.3729}.
\fi
\ifcsname pubref@BousseauSkein\endcsname
\bibitem{BousseauSkein} P. Bousseau,
\emph{Strong positivity for the skein algebras of the $4$-punctured sphere and the $1$-punctured torus},
Comm. Math. Phys. 398 (2023), 1--58,
\href{https://doi.org/10.1007/s00220-022-04512-9}{doi:10.1007/s00220-022-04512-9}.
\fi
\ifcsname pubref@EOR\endcsname
\bibitem{EOR} P. Etingof, A. Oblomkov, E. Rains,
\emph{Generalized double affine Hecke algebras of rank 1 and quantized del Pezzo surfaces},
Adv. Math. 212 (2007), 749--796,
\href{https://doi.org/10.1016/j.aim.2006.11.008}{doi:10.1016/j.aim.2006.11.008}.
\fi
\ifcsname pubref@CMR\endcsname
\bibitem{CMR} L. Chekhov, M. Mazzocco, V. Rubtsov,
\emph{Quantised Painlev\'e monodromy manifolds, Sklyanin and Calabi--Yau algebras},
Adv. Math. 376 (2021), 107442,
\href{https://doi.org/10.1016/j.aim.2020.107442}{doi:10.1016/j.aim.2020.107442},
\href{https://arxiv.org/abs/1905.02772}{arXiv:1905.02772}.
原著v2のDefinition 1.1、式(1.5)、Theorem 1.5、Proposition 4.2を参照。
\fi
\ifcsname pubref@Rains\endcsname
\bibitem{Rains} E. M. Rains,
\emph{Filtered deformations of elliptic algebras},
\href{https://arxiv.org/abs/2107.13540v2}{arXiv:2107.13540v2}, 17 February 2022,
Propositions 2.2, 2.3, Example 2.1, Section 6末尾のConjecture 1とRemark。
\fi
\ifcsname pubref@LVdB\endcsname
\bibitem{LVdB} W. Lowen, M. Van den Bergh,
\emph{Hochschild cohomology of abelian categories and ringed spaces},
Adv. Math. 198 (2005), 172--221,
\href{https://arxiv.org/abs/math/0405227}{arXiv:math/0405227}。
\fi
\ifcsname pubref@LowenObject\endcsname
\bibitem{LowenObject} W. Lowen,
\emph{Hochschild cohomology, the characteristic morphism and derived deformations},
Compos. Math. 144 (2008), 1557--1580,
\href{https://arxiv.org/abs/0707.2602}{arXiv:0707.2602}。
\fi
\ifcsname pubref@Repo\endcsname
\bibitem{Repo} \emph{nc-cubic-moduli research repository},
\href{https://github.com/uedakazushi/nc-cubic-moduli}{GitHub main}。
本稿の依存する個別ファイル・ラベル・確認済み状態は研究ノートの入力表に固定する。
\fi
\ifcsname pubref@CompanionE6\endcsname
\bibitem{CompanionE6}\emph{三乗根量子三次曲面の正常形と形式変形},
日本語研究原稿、2026年11月18日、\PublicationVersion。
\href{https://github.com/uedakazushi/nc-cubic-e6/blob/main/papers/nc-cubic-e6.tex}{\texttt{nc-cubic-e6.tex}}。
\fi
\ifcsname pubref@CompanionFricke\endcsname
\bibitem{CompanionFricke}\emph{量子三次曲面のFricke幾何と周期},
日本語研究原稿、2026年11月18日、\PublicationVersion。
\href{https://github.com/uedakazushi/nc-cubic-e6/blob/main/papers/fricke-geometry.tex}{\texttt{fricke-geometry.tex}}。
\fi
\ifcsname pubref@CompanionBoundary\endcsname
\bibitem{CompanionBoundary}\emph{三次曲面のPell力学と境界特異点},
日本語研究原稿、2026年11月18日、\PublicationVersion。
\href{https://github.com/uedakazushi/nc-cubic-e6/blob/main/papers/cubic-boundary-prym.tex}{\texttt{cubic-boundary-prym.tex}}。
\fi
\ifcsname pubref@AOUCubic\endcsname
\bibitem{AOUCubic}
T. Abdelgadir, S. Okawa, K. Ueda,
\emph{Compact moduli of marked noncommutative cubic surfaces},
\href{https://arxiv.org/abs/2404.00175v1}{arXiv:2404.00175v1} (2024).
\S9, equations (9.1)--(9.10), Lemma 9.2, Theorem 9.4.
\fi
\ifcsname pubref@OUASSurfaces\endcsname
\bibitem{OUASSurfaces}
S. Okawa, K. Ueda,
\emph{Marked Artin--Schelter surfaces of del Pezzo types},
\href{https://arxiv.org/abs/2508.04986v1}{arXiv:2508.04986v1} (2025).
Definition 1.1.
\fi
\ifcsname pubref@BoundaryDHKK\endcsname
\bibitem{BoundaryDHKK} G. Dimitrov, F. Haiden, L. Katzarkov, M. Kontsevich,
\emph{Dynamical systems and categories},
\href{https://arxiv.org/abs/1307.8418}{arXiv:1307.8418},
Definitions 2.1, 2.4, Proposition 2.2, Lemma 2.3, Theorem 2.6.
\fi
\ifcsname pubref@BoundaryGoldmanNeumann\endcsname
\bibitem{BoundaryGoldmanNeumann} W. M. Goldman, W. D. Neumann,
\emph{Homological action of the modular group on some cubic moduli spaces},
Math. Res. Lett. 12 (2005), 575--591,
\href{https://doi.org/10.4310/MRL.2005.v12.n4.a11}{doi:10.4310/MRL.2005.v12.n4.a11},
\href{https://arxiv.org/abs/math/0402039}{arXiv:math/0402039},
Theorems 1--2, Proposition 3.
\fi
\ifcsname pubref@BoundaryDyckerhoff\endcsname
\bibitem{BoundaryDyckerhoff} T. Dyckerhoff,
\emph{Compact generators in categories of matrix factorizations},
Duke Math. J. 159 (2011), 223--274,
\href{https://doi.org/10.1215/00127094-1415869}{doi:10.1215/00127094-1415869},
\href{https://arxiv.org/abs/0904.4713v5}{arXiv:0904.4713v5},
Theorem 5.9, Corollary 6.5.
\fi
\ifcsname pubref@BoundaryHuaKeller\endcsname
\bibitem{BoundaryHuaKeller} Z. Hua, B. Keller,
\emph{Cluster categories and rational curves},
Geom. Topol. 28 (2024), 2569--2634,
\href{https://doi.org/10.2140/gt.2024.28.2569}{doi:10.2140/gt.2024.28.2569},
\href{https://arxiv.org/abs/1810.00749v6}{arXiv:1810.00749v6},
Theorem 5.9, pp.36--37.
\fi
\ifcsname pubref@BoundaryMatherYau\endcsname
\bibitem{BoundaryMatherYau} J. N. Mather, S. S.-T. Yau,
\emph{Classification of isolated hypersurface singularities by their moduli algebras},
Invent. Math. 69 (1982), 243--251,
\href{https://doi.org/10.1007/BF01399504}{doi:10.1007/BF01399504}.
\fi
\ifcsname pubref@BoundaryMumford\endcsname
\bibitem{BoundaryMumford} D. Mumford,
\emph{Prym Varieties I}, in \emph{Contributions to Analysis},
Academic Press (1974), 325--350,
\href{https://doi.org/10.1016/B978-0-12-044850-0.50032-0}{doi:10.1016/B978-0-12-044850-0.50032-0},
\S\S1--3.
\fi
\ifcsname pubref@BoundaryMirandaPersson\endcsname
\bibitem{BoundaryMirandaPersson} R. Miranda, U. Persson,
\emph{On extremal rational elliptic surfaces},
Math. Z. 193 (1986), 537--558,
\href{https://doi.org/10.1007/BF01160474}{doi:10.1007/BF01160474},
Table 5.2.
\fi
\ifcsname pubref@BoundaryHowald\endcsname
\bibitem{BoundaryHowald} J. Howald,
\emph{Multiplier Ideals of Sufficiently General Polynomials},
\href{https://arxiv.org/abs/math/0303203}{arXiv:math/0303203},
Theorems 11--12, Corollary 13.
\fi
\ifcsname pubref@BoundaryCasals\endcsname
\bibitem{BoundaryCasals} R. Casals,
\emph{Lagrangian skeleta and plane curve singularities},
J. Fixed Point Theory Appl. 24 (2022), 34,
\href{https://doi.org/10.1007/s11784-022-00939-8}{doi:10.1007/s11784-022-00939-8},
Theorem 1.1, Corollary 1.2, \S2.1.
\fi
\ifcsname pubref@BoundaryCamara\endcsname
\bibitem{BoundaryCamara} L. M. C\^amara,
\emph{On the classification of quasi-homogeneous curves},
\href{https://arxiv.org/abs/1009.1664v1}{arXiv:1009.1664v1},
Theorem 1(iii), \S4.3.
\fi
\ifcsname pubref@FrickeFriedman\endcsname
\bibitem{FrickeFriedman}
R. Friedman, \emph{On the geometry of anticanonical pairs},
\href{https://arxiv.org/abs/1502.02560v2}{arXiv:1502.02560v2} (2016),
Proposition 3.12.
\fi
\ifcsname pubref@FrickeBoalchPaluba\endcsname
\bibitem{FrickeBoalchPaluba}
P. Boalch, R. Paluba, \emph{Symmetric cubic surfaces and $G_2$ character varieties},
J. Algebraic Geom. 25 (2016), 607--631,
\href{https://doi.org/10.1090/jag/668}{doi:10.1090/jag/668},
Lemma 6, Corollary 11, Proposition 18, Remark 19.
\fi
\ifcsname pubref@PymR13\endcsname
\bibitem{PymR13} B. Pym,
\emph{Quantum deformations of projective three-space},
Adv. Math. 281 (2015), 1216--1241,
\href{https://doi.org/10.1016/j.aim.2015.06.005}{doi:10.1016/j.aim.2015.06.005},
\href{https://arxiv.org/abs/1403.6444}{arXiv:1403.6444}, \S3.4, (9).
\fi
\ifcsname pubref@CompanionAmbient\endcsname
\bibitem{CompanionAmbient}
\emph{非可換三次曲面のambient射と$R(1,3)$成分},
\href{https://github.com/uedakazushi/nc-cubic-ambient/blob/main/papers/nc-cubic-ambient.tex}
{nc-cubic-ambient research manuscript},
\texttt{thm:r13-affine}, \texttt{prop:center}, \texttt{sec:surface-fiber}.
引用範囲は次数付き中心拡大の同型商と中心三次ペンシルであり、
AOUからの支配性の追加仮定は用いない。
\fi
\end{thebibliography}

\end{document}
